% !TEX program = pdflatex
\documentclass[11pt]{article}

\usepackage[T1]{fontenc}
\usepackage{lmodern}
\usepackage{microtype}
\usepackage[a4paper,margin=28mm,headheight=14pt]{geometry}
\usepackage{amsmath,amssymb,amsthm,mathtools}
\usepackage{enumitem}
\usepackage{array,booktabs,longtable}
\usepackage{xcolor}
\usepackage{xurl}
\usepackage{hyperref}
\usepackage{fancyhdr}

\definecolor{linkblue}{RGB}{24,78,119}
\definecolor{citegreen}{RGB}{36,104,71}
\hypersetup{
  colorlinks=true,
  linkcolor=linkblue,
  citecolor=citegreen,
  urlcolor=linkblue,
  pdftitle={A Dirichlet Polya Inequality for Three-Dimensional Spherical Shells},
  pdfsubject={Purely analytic main proof manuscript},
  pdfkeywords={Polya conjecture, Dirichlet Laplacian, spherical shell, eigenvalue counting}
}

\pagestyle{fancy}
\fancyhf{}
\fancyhead[L]{\small Dirichlet P\'olya inequality for spherical shells}
\fancyhead[R]{\small Purely analytic proof}
\fancyfoot[C]{\thepage}
\fancypagestyle{plain}{%
  \fancyhf{}%
  \fancyhead[L]{\small Dirichlet P\'olya inequality for spherical shells}%
  \fancyhead[R]{\small Purely analytic proof}%
  \fancyfoot[C]{\thepage}%
}
\fancypagestyle{titlepage}{%
  \fancyhf{}%
  \fancyfoot[C]{\thepage}%
  \renewcommand{\headrulewidth}{0pt}%
}

\numberwithin{equation}{section}
\allowdisplaybreaks[2]
\setlist{itemsep=0.25em,topsep=0.45em}
\setlength{\emergencystretch}{3em}
\newtheorem{theorem}{Theorem}[section]
\newtheorem{proposition}[theorem]{Proposition}
\newtheorem{lemma}[theorem]{Lemma}
\newtheorem{corollary}[theorem]{Corollary}
\theoremstyle{definition}
\newtheorem{definition}[theorem]{Definition}
\theoremstyle{remark}
\newtheorem{remark}[theorem]{Remark}

\newcommand{\R}{\mathbb R}
\newcommand{\N}{\mathbb N}
\newcommand{\Nzero}{\mathbb N_0}
\DeclareMathOperator{\arsinh}{arsinh}

\title{A Dirichlet P\'olya Inequality for\\Three-Dimensional Spherical Shells\\[0.4em]
\large Purely analytic main proof manuscript}
\author{}
\date{16 July 2026}

\begin{document}

\maketitle
\thispagestyle{titlepage}

\begin{quote}
\small
\textbf{Scope and status.}
This manuscript proves a theorem for the class of genuine three-dimensional
spherical shells.  It does not claim the general P\'olya conjecture,
literature novelty, or priority.  All shell-specific analytic reductions,
regional arguments, seam assignments, exact constants, and all finite
exact-rational decisions are stated in this paper and its accompanying
purely analytic supplement.  Executable replay sources and interval
certificates are retained only as optional regression checks and are not
premises of the proof.  The imported background and analytic results
comprise explicitly identified standard Sturm--Liouville and Bessel facts and
published phase, Bessel-zero, and one-dimensional floor-sum inequalities,
each quoted with its precise scope.
\end{quote}

\begin{abstract}
For \(0<r<R\), let
\[
  A_{r,R}=\{x\in\R^3:r<|x|<R\}.
\]
We prove, with the strict Dirichlet counting convention, that
\[
  N_D(A_{r,R},\Lambda)
  \le L_3|A_{r,R}|\Lambda^{3/2},
  \qquad
  L_3=\frac1{6\pi^2},
  \qquad \Lambda\ge0.
\]
The proof is purely analytic.  It separates variables, converts the radial
spectrum into a strict multiplicity-weighted phase count, and combines
analytic endpoint estimates, finite angular and radial mode inventories,
hand-checkable exact-rational ledgers, a retained-remainder estimate on the
final compact window, and an analytic curvature propagation for the frequency
tail.  No floating-point or interval-arithmetic conclusion is used as a proof
premise.  We also give an elementary proof that no genuine spherical
shell tiles \(\R^3\) by congruent rigid-motion copies with disjoint interiors,
under either exact or almost-everywhere coverage.
\end{abstract}

\tableofcontents
\clearpage

\section{Main results and conventions}

The eigenvalues of the Dirichlet Laplacian on a bounded domain \(\Omega\) are
listed with multiplicity as
\[
  0<\lambda_1^D(\Omega)\le\lambda_2^D(\Omega)\le\cdots.
\]
Throughout we use the strict counting function
\begin{equation}
  N_D(\Omega,\Lambda)
  :=\#\{j:\lambda_j^D(\Omega)<\Lambda\}.
  \label{eq:strict-count}
\end{equation}
This convention matters at every spectral wall.  We write
\(\N=\{1,2,3,\ldots\}\) and \(\Nzero=\{0,1,2,\ldots\}\).

\begin{theorem}[Spectral inequality]
\label{thm:spectral}
For every \(0<r<R\) and every \(\Lambda\ge0\),
\begin{equation}
  \boxed{
  N_D(A_{r,R},\Lambda)
  \le \frac{2}{9\pi}(R^3-r^3)\Lambda^{3/2}
  =L_3|A_{r,R}|\Lambda^{3/2},
  \qquad L_3=\frac1{6\pi^2}.}
  \label{eq:main-theorem}
\end{equation}
\end{theorem}

\begin{theorem}[Non-tiling]
\label{thm:nontiling}
For every \(0<r<R\), no family of congruent rigid-motion copies of
\(A_{r,R}\) has pairwise-disjoint interiors and covers \(\R^3\) exactly or
up to a three-dimensional Lebesgue-null set.  The same conclusion holds when
coverage is formulated using closed shells.
\end{theorem}

The two theorems concern literally the same complete shell family.
Theorem~\ref{thm:nontiling} is logically independent of the spectral proof.

% Self-contained manuscript fragment: separated spectrum, phase reduction,
% weighted lattice reduction, and the small-hole endpoint.
% Assumes amsmath, amssymb, amsthm, mathtools and theorem environments
% theorem, proposition, lemma.

\section{Separated spectrum, phase reduction, and the small-hole endpoint}
\label{sec:sc-spectrum-small-hole}

Throughout this section
\[
  A_{\rho,1}=\{x\in\mathbb R^3:\rho<|x|<1\},
  \qquad 0<\rho<1,
\]
and the Dirichlet counting function is strict:
\[
  N_D(A_{\rho,1},K^2)
  :=\#\{j:\lambda_j^D(A_{\rho,1})<K^2\}.
\]
For \(x>0\) we use the strict positive-integer count
\begin{equation}
  [x]_<:=\#\{n\in\mathbb N:n<x\}
  =\max\{0,\lceil x\rceil-1\}.
  \label{eq:sc-strict-bracket}
\end{equation}
Thus \([m]_< =m-1\) for \(m\in\mathbb N\).  This convention is retained
at every eigenfrequency, phase level, and half-integer angular wall.

\subsection{The exact separated spectrum}

Let \(H_\rho=-\Delta^D_{A_{\rho,1}}\) be the Friedrichs operator associated
with
\[
  \mathfrak q[f]=\int_{A_{\rho,1}}|\nabla f|^2\,dx,
  \qquad D(\mathfrak q)=H_0^1(A_{\rho,1}).
\]
Choose an orthonormal spherical-harmonic basis satisfying
\[
  -\Delta_{\mathbb S^2}Y_{\ell m}
  =\ell(\ell+1)Y_{\ell m},
  \qquad \ell\in\mathbb N_0,
  \quad -\ell\le m\le\ell.
\]

\begin{proposition}[Unitary decomposition and completeness]
\label{prop:sc-separated-spectrum}
For
\[
  L_\ell=-\frac{d^2}{dr^2}+\frac{\ell(\ell+1)}{r^2},
  \qquad
  D(L_\ell)=H^2(\rho,1)\cap H_0^1(\rho,1),
\]
one has the unitary equivalence
\begin{equation}
  H_\rho\simeq
  \bigoplus_{\ell=0}^{\infty}
  \bigoplus_{m=-\ell}^{\ell}L_\ell.
  \label{eq:sc-operator-sum}
\end{equation}
The exact form domain on the right is
\begin{equation}
  \left\{(u_{\ell m}):
  \begin{array}{l}
  u_{\ell m}\in H_0^1(\rho,1),\\
  \displaystyle
  \sum_{\ell,m}\left(\|u_{\ell m}\|_2^2
  +\mathfrak a_\ell[u_{\ell m}]\right)<\infty
  \end{array}\right\},
  \label{eq:sc-global-form-domain}
\end{equation}
where
\[
  \mathfrak a_\ell[u]
  =\int_\rho^1\left(|u'|^2
  +\frac{\ell(\ell+1)}{r^2}|u|^2\right)dr.
\]
The exact operator domain is
\begin{equation}
  \left\{(u_{\ell m}):
  \begin{array}{l}
  u_{\ell m}\in D(L_\ell),\\
  \displaystyle
  \sum_{\ell,m}\left(\|u_{\ell m}\|_2^2
  +\|L_\ell u_{\ell m}\|_2^2\right)<\infty
  \end{array}\right\}.
  \label{eq:sc-global-operator-domain}
\end{equation}
Each \(L_\ell\) has positive, simple eigenvalues tending to infinity and a
complete orthonormal eigenbasis.  The full operator has compact resolvent.
\end{proposition}

\begin{proof}
Write \(x=r\omega\), so \(dx=r^2\,dr\,d\omega\), and set
\[
  R_{\ell m}(r)=\int_{\mathbb S^2}
  f(r,\omega)\overline{Y_{\ell m}(\omega)}\,d\omega,
  \qquad u_{\ell m}(r)=rR_{\ell m}(r).
\]
Fubini, Tonelli, and Parseval on \(\mathbb S^2\) give
\[
  \|f\|_{L^2(A_{\rho,1})}^2
  =\sum_{\ell,m}\int_\rho^1|R_{\ell m}(r)|^2r^2\,dr
  =\sum_{\ell,m}\int_\rho^1|u_{\ell m}(r)|^2\,dr.
\]
Conversely, if \((u_{\ell m})\) is square summable, the finite sums
\[
  f_L(r,\omega)=\frac1r
  \sum_{\ell\le L}\sum_{m=-\ell}^{\ell}
  u_{\ell m}(r)Y_{\ell m}(\omega)
\]
are Cauchy in \(L^2(A_{\rho,1})\), and their limit maps back to the given
family.  Hence the coefficient map is unitary onto the Hilbert direct sum.
Because \(r,1/r\in W^{1,\infty}(\rho,1)\), multiplication by either function
is bounded on \(H^1(\rho,1)\) and preserves the two zero endpoint traces.

For a finite smooth harmonic expansion, polar coordinates and angular
orthogonality yield
\[
  \mathfrak q[f]=\sum_{\ell,m}\int_\rho^1
  \left(r^2|R_{\ell m}'|^2
  +\ell(\ell+1)|R_{\ell m}|^2\right)dr.
\]
If \(u=rR\), then
\[
  r^2|R'|^2
  =\left|u'-\frac ur\right|^2
  =|u'|^2-\frac{(|u|^2)'}r+\frac{|u|^2}{r^2}.
\]
Since
\[
  \frac{(|u|^2)'}r
  =\left(\frac{|u|^2}{r}\right)'
  +\frac{|u|^2}{r^2}
\]
and \(u(\rho)=u(1)=0\), integration gives
\[
  \int_\rho^1r^2|R'|^2\,dr
  =\int_\rho^1|u'|^2\,dr.
\]
The angular term becomes
\(\ell(\ell+1)|u|^2/r^2\).  Closing finite harmonic expansions in the
form norm proves \eqref{eq:sc-global-form-domain}: one inclusion follows
from componentwise form convergence and Fatou's lemma, while the reverse
inclusion follows by first truncating the square-summable family in
\((\ell,m)\) and then approximating its finitely many components by
\(C_c^\infty(\rho,1)\).  The uniqueness theorem for operators associated
with closed semibounded forms gives \eqref{eq:sc-operator-sum}.

For fixed \(\ell\), the potential \(\ell(\ell+1)/r^2\) is continuous and
bounded on \([\rho,1]\).  The regular one-dimensional Dirichlet form
therefore has associated operator \(L_\ell\) with the displayed
\(H^2\cap H_0^1\) domain.  The standard domain formula for an orthogonal
operator sum gives \eqref{eq:sc-global-operator-domain}.

The embedding \(H_0^1(\rho,1)\hookrightarrow L^2(\rho,1)\) is compact, so
every block has compact resolvent.  To pass to the infinite direct sum, note
that \(r\le1\) implies
\[
  L_\ell\ge\ell(\ell+1),
  \qquad
  \|(L_\ell+1)^{-1}\|
  \le\frac1{1+\ell(\ell+1)}.
\]
The resolvent of the full operator is consequently the operator-norm limit
of its finite-\(\ell\) truncations, each of which is compact.  Hence the full
resolvent is compact.

Regular Dirichlet Sturm--Liouville theory gives a complete sequence of
simple eigenvalues in every block.  They are strictly positive because
\[
  \mathfrak a_\ell[u]\ge\int_\rho^1|u'|^2\,dr
  \ge\frac{\pi^2}{(1-\rho)^2}\|u\|_2^2.
\]
Simplicity also follows directly: two solutions of the same second-order
equation that vanish at \(r=\rho\) have proportional initial derivatives
and hence are proportional by uniqueness.  If \(u_{\ell,n}\) is a complete
radial eigenbasis, then
\[
  \frac{u_{\ell,n}(r)}rY_{\ell m}(\omega)
\]
over all \((\ell,m,n)\) is exactly the inverse image of a complete
orthonormal coordinate basis of the Hilbert direct sum.  This proves global
completeness.
\end{proof}

\begin{proposition}[Bessel determinant and strict phase enumeration]
\label{prop:sc-bessel-phase-count}
Put \(\nu=\ell+\tfrac12\).  The positive eigenfrequencies in the
\(\ell\)-th radial block are precisely the positive zeros of
\begin{equation}
  D_{\nu,\rho}(k)
  :=J_\nu(\rho k)Y_\nu(k)-Y_\nu(\rho k)J_\nu(k).
  \label{eq:sc-bessel-determinant}
\end{equation}
Every such zero is simple.  With the continuous phase branch
\begin{equation}
  J_\nu(x)+iY_\nu(x)=M_\nu(x)e^{i\theta_\nu(x)},
  \qquad M_\nu(x)>0,
  \qquad \theta_\nu(0+)=-\frac\pi2,
  \label{eq:sc-phase-branch}
\end{equation}
and
\[
  \Psi_{\nu,\rho}(K)=\theta_\nu(K)-\theta_\nu(\rho K),
  \qquad
  \gamma_{\rho,K}(\nu)=\frac{\Psi_{\nu,\rho}(K)}\pi,
\]
the exact strict count is
\begin{equation}
  N_D(A_{\rho,1},K^2)
  =\sum_{\ell=0}^{\infty}(2\ell+1)
  [\gamma_{\rho,K}(\ell+\tfrac12)]_<.
  \label{eq:sc-exact-phase-count}
\end{equation}
The sum is finite.  If \(K\) is a radial eigenfrequency in one or several
angular channels, the root at \(K\) is excluded in each such channel and the
excluded angular multiplicities add.
\end{proposition}

\begin{proof}
For \(k>0\), the physical radial equation has the general solution
\[
  R(r)=r^{-1/2}
  \bigl(c_1J_\nu(kr)+c_2Y_\nu(kr)\bigr).
\]
The two Dirichlet conditions admit a nonzero coefficient vector exactly when
\[
  \det\begin{pmatrix}
  J_\nu(\rho k)&Y_\nu(\rho k)\\
  J_\nu(k)&Y_\nu(k)
  \end{pmatrix}=D_{\nu,\rho}(k)=0.
\]
The Wronskian
\begin{equation}
  J_\nu(x)Y_\nu'(x)-J_\nu'(x)Y_\nu(x)=\frac2{\pi x}
  \label{eq:sc-bessel-wronskian}
\end{equation}
shows that the two entries in either row cannot vanish simultaneously.
Thus the determinant criterion remains exact even if the \(J_\nu\) values,
or the \(Y_\nu\) values, happen to vanish at both endpoints.

For root simplicity define the left-normalized solution
\begin{equation}
  s_\ell(r,k)=\frac{\pi\sqrt\rho}{2}\sqrt r
  \left[J_\nu(\rho k)Y_\nu(kr)
  -Y_\nu(\rho k)J_\nu(kr)\right].
  \label{eq:sc-left-normalized}
\end{equation}
Equation \eqref{eq:sc-bessel-wronskian} gives
\[
  s_\ell(\rho,k)=0,
  \qquad \partial_rs_\ell(\rho,k)=1,
  \qquad
  s_\ell(1,k)=\frac{\pi\sqrt\rho}{2}D_{\nu,\rho}(k).
\]
Let \(v=\partial_ks_\ell\).  Differentiating
\((L_\ell-k^2)s_\ell=0\) gives
\[
  (L_\ell-k^2)v=2ks_\ell,
  \qquad v(\rho,k)=\partial_rv(\rho,k)=0.
\]
At a positive determinant root, Lagrange's identity yields
\begin{equation}
  \partial_ks_\ell(1,k)\,\partial_rs_\ell(1,k)
  =2k\int_\rho^1s_\ell(r,k)^2\,dr>0.
  \label{eq:sc-root-simplicity}
\end{equation}
Thus \(\partial_ks_\ell(1,k)\ne0\), and the zero of the determinant is
simple.

At zero energy the left-normalized solution is
\[
  s_\ell(r,0)=
  \frac{\rho^{-\ell}r^{\ell+1}
  -\rho^{\ell+1}r^{-\ell}}{2\ell+1},
\]
so
\[
  s_\ell(1,0)=
  \frac{\rho^{-\ell}-\rho^{\ell+1}}{2\ell+1}>0.
\]
Equivalently, the standard small-argument expansions give
\begin{equation}
  \lim_{k\downarrow0}D_{\nu,\rho}(k)
  =\frac{\rho^{-\nu}-\rho^\nu}{\pi\nu}>0.
  \label{eq:sc-zero-determinant-limit}
\end{equation}
Hence the limiting phase level zero is not an eigenfrequency.

From \eqref{eq:sc-phase-branch},
\[
  J_\nu=M_\nu\cos\theta_\nu,
  \qquad Y_\nu=M_\nu\sin\theta_\nu,
\]
and direct substitution gives the sign-sensitive factorization
\begin{equation}
  D_{\nu,\rho}(k)
  =M_\nu(\rho k)M_\nu(k)
  \sin\Psi_{\nu,\rho}(k).
  \label{eq:sc-determinant-factorization}
\end{equation}
The Wronskian also gives
\[
  \theta_\nu'(x)=\frac2{\pi xM_\nu(x)^2}>0.
\]
Nicholson's integral representation
\[
  M_\nu(x)^2=\frac8{\pi^2}\int_0^\infty
  \cosh(2\nu t)K_0(2x\sinh t)\,dt
\]
and the strict decrease of \(K_0\) imply that \(M_\nu(x)^2\) is strictly
decreasing in \(x\).  Therefore
\begin{equation}
  \Psi_{\nu,\rho}'(k)
  =\frac2{\pi k}
  \left(M_\nu(k)^{-2}-M_\nu(\rho k)^{-2}\right)>0.
  \label{eq:sc-phase-monotonicity}
\end{equation}
The branch condition gives \(\Psi_{\nu,\rho}(0+)=0\), while the standard
large-argument Bessel expansion gives
\[
  \Psi_{\nu,\rho}(k)=(1-\rho)k+O(k^{-1})\longrightarrow\infty
\]
for fixed \(\nu\) and \(\rho\).  Hence the positive determinant roots are
uniquely indexed by
\[
  \Psi_{\ell+1/2,\rho}(k_{\ell,n})=n\pi,
  \qquad n=1,2,\ldots.
\]
Strict monotonicity gives
\[
  k_{\ell,n}<K
  \quad\Longleftrightarrow\quad
  n\pi<\Psi_{\ell+1/2,\rho}(K).
\]
The direct-sum decomposition supplies angular multiplicity \(2\ell+1\).
If the same numerical eigenvalue occurs at distinct values of \(\ell\), the
corresponding spherical-harmonic sectors are orthogonal and their
multiplicities add.  Finally,
\(\lambda_{\ell,n}\ge\ell(\ell+1)\), so only finitely many angular channels
can contribute below \(K^2\).  These observations prove
\eqref{eq:sc-exact-phase-count}, including every strict spectral wall.
\end{proof}

The only standard inputs in the preceding proof are completeness of the
spherical harmonics, regular one-dimensional Dirichlet Sturm--Liouville
theory, the Bessel equation, the Wronskian, Nicholson's representation, and
the usual small- and large-argument Bessel formulas; see \cite{DLMF}.

\subsection{Uniform phase enclosures and the strict lattice proxy}

For \(a>0\), define
\begin{equation}
  G_a(z)=
  \begin{cases}
  \displaystyle\frac1\pi\left(\sqrt{a^2-z^2}
  -z\arccos\frac za\right),&0\le z\le a,\\[1mm]
  0,&z>a.
  \end{cases}
  \label{eq:sc-model-action}
\end{equation}
For \(0\le z<a\), put
\[
  H_a(z)=\frac{3a^2+2z^2}{24\pi(a^2-z^2)^{3/2}},
\]
and define
\begin{equation}
  \mathcal F_a(z)=
  \begin{cases}
  \max\{G_a(z)-H_a(z),-1/4\},&0\le z<a,\\
  -1/4,&z\ge a.
  \end{cases}
  \label{eq:sc-auxiliary-F}
\end{equation}

We use the following published Bessel-phase result, stated here with every
hypothesis needed in the sequel.

\begin{lemma}[External uniform Bessel-phase enclosure]
\label{lem:sc-external-phase}
With the branch \eqref{eq:sc-phase-branch}, for every \(a>0\) and real
\(z\ge0\),
\begin{equation}
  \mathcal F_a(z)-\frac14
  <\frac{\theta_z(a)}\pi
  <G_a(z)-\frac14.
  \label{eq:sc-individual-phase-enclosure}
\end{equation}
If \(0<\mu<K\), \(0\le z\le K\), and
\[
  \Gamma_{K,\mu}(z)
  :=\frac{\theta_z(K)-\theta_z(\mu)}\pi,
\]
then
\begin{equation}
  0<\Gamma_{K,\mu}(z)
  <G_K(z)-\mathcal F_\mu(z)
  \le G_K(z)+\frac14.
  \label{eq:sc-global-difference-enclosure}
\end{equation}
If \(0\le z\le\mu\), including \(z=\mu\), then also
\begin{equation}
  G_K(z)-G_\mu(z)
  <\Gamma_{K,\mu}(z)
  <G_K(z)-G_\mu(z)+\frac14.
  \label{eq:sc-action-difference-enclosure}
\end{equation}
No value of the divergent expression \(H_\mu(\mu)\) is used in the last
statement.
\end{lemma}

Lemma~\ref{lem:sc-external-phase} is Theorem~5.1 and Lemma~5.2 of
\cite{FLPSannuli}, based on the real-order phase enclosures of
\cite{FLPSphase}.  The upper bound in
\eqref{eq:sc-global-difference-enclosure} follows immediately by subtracting
the lower bound for \(\theta_z(\mu)\) from the upper bound for
\(\theta_z(K)\) in \eqref{eq:sc-individual-phase-enclosure}; the correlated
two-sided estimate \eqref{eq:sc-action-difference-enclosure} is the sharper
part of the annulus result.  Both theorems are stated for real order, so the
half-integers \(z=\ell+\tfrac12\) require no interpolation.

Fix now \(K>0\), set \(\mu=\rho K\), and define
\begin{equation}
  A_{\rho,K}(z)=G_K(z)-G_\mu(z).
  \label{eq:sc-shell-action}
\end{equation}
For \(0\le z\le K\), let
\begin{equation}
  s_\mu(z)=
  \begin{cases}
  \min\{H_\mu(z),1/4\},&0\le z<\mu,\\
  1/4,&\mu\le z\le K,
  \end{cases}
  \qquad
  U_{\rho,K}(z)=A_{\rho,K}(z)+s_\mu(z).
  \label{eq:sc-optimized-proxy}
\end{equation}
The second line, rather than \(H_\mu\), is always used at \(z=\mu\).

\begin{proposition}[Strict phase-to-lattice reduction]
\label{prop:sc-phase-lattice}
For \(0<\rho<1\) and \(K>0\),
\begin{equation}
  \gamma_{\rho,K}(z)<U_{\rho,K}(z)
  \le A_{\rho,K}(z)+\frac14,
  \qquad 0\le z\le K.
  \label{eq:sc-pointwise-phase-proxy}
\end{equation}
Because every \(G_a\) is zero-extended, \(A_{\rho,K}\) is the piecewise
action \(G_K-G_{\rho K}\) below the inner interface and \(G_K\) at and
above it.  Consequently,
\begin{align}
  N_D(A_{\rho,1},K^2)
  &\le
  \sum_{\nu_\ell<K}2\nu_\ell
  [U_{\rho,K}(\nu_\ell)]_<
  \label{eq:sc-optimized-spectral-proxy}\\
  &\le
  S_{\rho,K}:=
  \sum_{\nu_\ell<K}2\nu_\ell
  \left\lfloor A_{\rho,K}(\nu_\ell)+\frac14\right\rfloor,
  \qquad \nu_\ell=\ell+\frac12.
  \label{eq:sc-coarse-spectral-proxy}
\end{align}
The cutoff is strict: an order \(\nu_\ell=K\) is inactive.  At a true phase
wall the radial count is one less than the ordinary floor; at a proxy wall
the ordinary floor in \eqref{eq:sc-coarse-spectral-proxy} is retained.
Moreover,
\begin{equation}
  \int_0^\infty2zA_{\rho,K}(z)\,dz
  =\frac{2}{9\pi}(1-\rho^3)K^3.
  \label{eq:sc-exact-action-integral}
\end{equation}
\end{proposition}

\begin{proof}
For \(0\le z<\mu\), the first upper bound in
\eqref{eq:sc-global-difference-enclosure} reads
\[
  \gamma_{\rho,K}(z)
  <A_{\rho,K}(z)+G_\mu(z)-\mathcal F_\mu(z).
\]
From \eqref{eq:sc-auxiliary-F},
\[
  G_\mu(z)-\mathcal F_\mu(z)
  =\min\{H_\mu(z),G_\mu(z)+1/4\}.
\]
The upper half of \eqref{eq:sc-action-difference-enclosure} also gives
\(\gamma_{\rho,K}(z)<A_{\rho,K}(z)+1/4\).  Taking the better of these two
strict upper bounds, and using \(G_\mu(z)+1/4\ge1/4\), gives
\[
  \gamma_{\rho,K}(z)
  <A_{\rho,K}(z)+\min\{H_\mu(z),1/4\}
  =U_{\rho,K}(z).
\]
For \(\mu\le z\le K\), one has \(G_\mu(z)=0\), and
\eqref{eq:sc-global-difference-enclosure} gives
\[
  \gamma_{\rho,K}(z)<G_K(z)+1/4=U_{\rho,K}(z).
\]

Orders \(z\ge K\) do not contribute.  Indeed,
\eqref{eq:sc-individual-phase-enclosure} gives
\[
  \frac{\theta_z(K)}\pi<-\frac14,
  \qquad
  \frac{\theta_z(\mu)}\pi>-\frac12,
\]
because \(G_K(z)=0\) and \(\mathcal F_\mu(z)=-1/4\).  Together with phase
monotonicity this yields
\[
  0<\gamma_{\rho,K}(z)<\frac14<1,
\]
so there is no positive phase level below \(K\).  This includes \(z=K\).

For positive \(x<y\), every positive integer strictly below \(x\) is also
strictly below \(y\), so \([x]_<\le[y]_<\le\lfloor y\rfloor\).  Applying
this observation to the exact phase count
\eqref{eq:sc-exact-phase-count}, and then using
\(s_\mu\le1/4\), proves
\eqref{eq:sc-optimized-spectral-proxy}--
\eqref{eq:sc-coarse-spectral-proxy} without changing either wall
convention.

Finally, scaling \(z=ax\) gives
\[
  \int_0^a2zG_a(z)\,dz
  =\frac{2a^3}{\pi}
  \left(\int_0^1x\sqrt{1-x^2}\,dx
  -\int_0^1x^2\arccos x\,dx\right).
\]
The two elementary integrals are \(1/3\) and \(2/9\), respectively; hence
the last expression is \(2a^3/(9\pi)\).  Subtracting the result with
\(a=\mu=\rho K\) from that with \(a=K\) proves
\eqref{eq:sc-exact-action-integral}.
\end{proof}

The necessity of retaining integerization can be displayed exactly.  Put
\[
  m_K=\max\{0,\lceil K-\tfrac12\rceil\},
  \qquad
  r_\mu=\max\{0,\lceil\mu-\tfrac12\rceil\}.
\]
Thus the active orders are \(\nu_0,\ldots,\nu_{m_K-1}\), and precisely the
first \(r_\mu\) of them lie strictly below \(\mu\).  Define
\[
  E_0(a)=\sum_{\nu_\ell<a}2\nu_\ell G_a(\nu_\ell)
  -\frac{2a^3}{9\pi},
  \qquad
  E_{\mathrm{mp}}(\rho,K)=E_0(K)-E_0(\mu),
\]
\[
  E_{\mathrm{ph}}
  =\sum_{\ell=0}^{r_\mu-1}(2\ell+1)
  \min\{H_\mu(\nu_\ell),1/4\}
  +\frac{m_K^2-r_\mu^2}{4}.
\]
For \(x>0\), define the strict remainder
\[
  \langle x\rangle_<:=x-[x]_<\in(0,1],
\]
so \(\langle m\rangle_<=1\) at positive integer \(m\), and set
\[
  E_{\mathrm{frac}}
  =\sum_{\nu_\ell<K}2\nu_\ell
  \langle U_{\rho,K}(\nu_\ell)\rangle_<.
\]
Expanding \(U=A+s_\mu\) term by term and using
\eqref{eq:sc-exact-action-integral} gives the exact identity
\begin{equation}
  \sum_{\nu_\ell<K}2\nu_\ell[U_{\rho,K}(\nu_\ell)]_<
  -\frac{2}{9\pi}(1-\rho^3)K^3
  =E_{\mathrm{mp}}+E_{\mathrm{ph}}-E_{\mathrm{frac}}.
  \label{eq:sc-fractional-budget}
\end{equation}
In particular, the blanket quarter-slack cost is exactly
\[
  \frac14\sum_{\ell=0}^{m_K-1}(2\ell+1)=\frac{m_K^2}{4};
\]
it cannot be discarded merely because it is of lower polynomial order for
fixed \(\rho\).

\subsection{Multiplicity reduction to shifted one-dimensional tails}

We record the two external floor-sum inequalities used in the small-hole
argument.  For integers \(a<b\), write
\begin{equation}
  \mathbf T(h;a,b)
  :=\frac12\lfloor h(a)\rfloor
  +\sum_{j=a+1}^{b-1}\lfloor h(j)\rfloor
  +\frac12\lfloor h(b)\rfloor.
  \label{eq:sc-trapezoidal-floor}
\end{equation}

\begin{lemma}[External trapezoidal floor bounds]
\label{lem:sc-external-floor-bounds}
The following statements hold.
\begin{enumerate}
  \item If \(h\) is nonnegative and concave on an integer interval
  \([a,b]\), then
  \begin{equation}
    \mathbf T(h;a,b)\le\int_a^bh(t)\,dt.
    \label{eq:sc-concave-floor-bound}
  \end{equation}
  \item Suppose \(h\) is nonnegative, decreasing, concave, and
  \(c\)-Lipschitz on an
  integer interval \([a,b]\), where \(0<c<1\).  If an integer
  \(p\in[a,b)\) is the last point of the first ordinary-floor shelf, so that
  \[
    \lfloor h(a)\rfloor=\lfloor h(p)\rfloor
    >\lfloor h(p+1)\rfloor,
  \]
  then
  \begin{equation}
    \mathbf T(h;a,b)
    \leq\int_a^b h(t)\,dt-
          \frac{1-c}{2}(b-p).
    \label{eq:sc-improved-concave-floor-bound}
  \end{equation}
  If no such drop occurs, set \(p=b\); the basic bound in item~1 gives
  the same displayed formula with zero correction.  This last no-drop case
  is a consequence of the basic proposition, not an invocation of the
  improved theorem at \(p=b\).
  \item If \(g\) is nonnegative, decreasing, convex, and
  \(1/2\)-Lipschitz on an integer interval \([a,b]\), if \(g(b)=0\), and if
  \(g\) is \(1/3\)-Lipschitz on \([t_*,b]\) for some
  \(t_*\in[a,b]\), then
  \begin{equation}
    \mathbf T(g+1/4;a,b)
    \le\int_a^bg(t)\,dt-\frac14\lfloor g(t_*)\rfloor.
    \label{eq:sc-improved-convex-floor-bound}
  \end{equation}
  \item If \(g\) is nonnegative, decreasing, convex, and
  \(1/2\)-Lipschitz on \([0,b]\), where \(b>0\) is arbitrary and \(g(b)=0\),
  then
  \begin{equation}
    \left\lfloor g(0)+\frac14\right\rfloor
    +2\sum_{j=1}^{\lfloor b\rfloor}
    \left\lfloor g(j)+\frac14\right\rfloor
    \le2\int_0^bg(t)\,dt.
    \label{eq:sc-convex-shifted-tail-bound}
  \end{equation}
\end{enumerate}
\end{lemma}

The first assertion is Proposition~3.1 of \cite{FLPSannuli}; the
first-shelf improvement is its Theorem~1.4; and the third assertion is its
Theorem~3.4.  The fourth is the translated convex shifted-floor theorem of
\cite{FLPSballs}.  These published one-dimensional inequalities are the only
non-special-function external inputs in what follows.

For \(r\in\mathbb N_0\), put \(x_r=r+\tfrac12\) and define the coarse shifted
tail
\begin{equation}
  \mathcal T_r=
  \left\lfloor A_{\rho,K}(x_r)+\frac14\right\rfloor
  +2\sum_{j\ge1}
  \left\lfloor A_{\rho,K}(x_r+j)+\frac14\right\rfloor.
  \label{eq:sc-shifted-tail}
\end{equation}
All but finitely many terms vanish.

\begin{lemma}[Weighted-lattice scaffold]
\label{lem:sc-weighted-scaffold}
One has
\begin{equation}
  S_{\rho,K}=\sum_{r\ge0}\mathcal T_r.
  \label{eq:sc-tail-decomposition}
\end{equation}
If, for every \(r\) with \(x_r<K\),
\begin{equation}
  \mathcal T_r\le2\int_{x_r}^{K}A_{\rho,K}(z)\,dz,
  \label{eq:sc-tail-target}
\end{equation}
then
\begin{equation}
  S_{\rho,K}\le
  \frac{2}{9\pi}(1-\rho^3)K^3.
  \label{eq:sc-scaffold-conclusion}
\end{equation}
\end{lemma}

\begin{proof}
For any finitely supported nonnegative integer sequence \(q_\ell\),
\[
  \sum_{\ell\ge0}(2\ell+1)q_\ell
  =\sum_{r\ge0}\left(q_r+2\sum_{\ell>r}q_\ell\right).
\]
Indeed, \(q_\ell\) appears once with coefficient one when \(r=\ell\) and
appears \(\ell\) times with coefficient two when \(0\le r<\ell\).  Taking
\[
  q_\ell=\left\lfloor
  A_{\rho,K}(\ell+\tfrac12)+\frac14\right\rfloor
\]
proves \eqref{eq:sc-tail-decomposition}; samples at or beyond \(K\) are
zero.

Assume \eqref{eq:sc-tail-target}.  Tonelli's theorem for the finite
nonnegative sum gives
\[
  S_{\rho,K}
  \le2\int_0^K f_3(z)A_{\rho,K}(z)\,dz,
  \qquad
  f_3(z)=\#\{r\in\mathbb N_0:r+\tfrac12\le z\}.
\]
Let \(F_3(z)=\int_0^zf_3(t)\,dt\).  If
\(z=m+\tfrac12+s\), where \(m\in\mathbb N_0\) and \(0\le s<1\), then direct
summation gives
\begin{equation}
  \frac{z^2}{2}-F_3(z)
  =\frac12\left(s-\frac12\right)^2\ge0;
  \label{eq:sc-step-primitive}
\end{equation}
for \(0\le z<1/2\) the same inequality is immediate because \(F_3(z)=0\).
The action \(A_{\rho,K}\) is absolutely continuous, decreasing, and vanishes
at \(K\).  Indeed, away from the inner interface,
\[
  A_{\rho,K}'(z)=
  \begin{cases}
  \displaystyle-\frac1\pi\left(
  \arccos\frac zK-\arccos\frac z\mu\right)<0,
  &0<z<\mu,\\[2mm]
  \displaystyle-\frac1\pi\arccos\frac zK<0,
  &\mu<z<K,
  \end{cases}
\]
and the two formulas have the same limiting value at \(z=\mu\).
Integration by parts and \eqref{eq:sc-step-primitive} therefore
give
\[
  2\int_0^Kf_3A_{\rho,K}
  =-2\int_0^KF_3A_{\rho,K}'
  \le-\int_0^Kz^2A_{\rho,K}'(z)\,dz
  =\int_0^K2zA_{\rho,K}(z)\,dz.
\]
Now apply \eqref{eq:sc-exact-action-integral}.
\end{proof}

\begin{lemma}[Tails at or above the inner interface]
\label{lem:sc-high-interface-tails}
If \(x_r\ge\mu=\rho K\) and \(x_r<K\), then
\[
  \mathcal T_r\le2\int_{x_r}^{K}A_{\rho,K}(z)\,dz.
\]
Equality \(x_r=\mu\) is included.
\end{lemma}

\begin{proof}
On \(z\ge x_r\ge\mu\), one has \(G_\mu(z)=0\), including at \(z=\mu\),
so \(A_{\rho,K}(z)=G_K(z)\).  Set \(b=K-x_r>0\) and
\(g(t)=G_K(x_r+t)\) on \([0,b]\).  Direct differentiation of
\eqref{eq:sc-model-action} gives
\[
  G_K'(z)=-\frac1\pi\arccos\frac zK\in[-1/2,0],
  \qquad
  G_K''(z)=\frac1{\pi\sqrt{K^2-z^2}}>0.
\]
Thus \(g\) is nonnegative, decreasing, convex, and \(1/2\)-Lipschitz, with
\(g(b)=G_K(K)=0\).  Terms of \eqref{eq:sc-shifted-tail} with \(j>b\) vanish;
if \(b\) is an integer, the terminal term is
\(\lfloor G_K(K)+1/4\rfloor=0\).  Therefore
\eqref{eq:sc-convex-shifted-tail-bound} is exactly the desired estimate.
\end{proof}

\subsection{The low-interface estimate in the small-hole sector}

Define
\begin{equation}
  \omega_0:=G_1(1/2)=\frac{\sqrt3}{2\pi}-\frac16,
  \qquad
  C_*:=\frac12+\frac{8\sqrt2}{15}.
  \label{eq:sc-small-hole-constants}
\end{equation}
Notice that \(0<\omega_0<1/2\).  The lower inequality is equivalent to
\(\pi<3\sqrt3\), which follows from \(\pi<4<3\sqrt3\); the upper inequality
also follows directly from \(0<G_1(1/2)<G_1(0)=1/\pi<1/2\).

\begin{lemma}[Low-interface shifted tails]
\label{lem:sc-low-interface-tails}
Assume
\begin{equation}
  0<\rho<\omega_0,
  \qquad
  K(\omega_0-\rho)\ge C_*.
  \label{eq:sc-low-interface-sector}
\end{equation}
If \(x_r=r+\tfrac12<\mu=\rho K\), then, in fact strictly,
\begin{equation}
  \mathcal T_r
  <2\int_{x_r}^{K}A_{\rho,K}(z)\,dz.
  \label{eq:sc-low-tail-conclusion}
\end{equation}
\end{lemma}

\begin{proof}
Fix a low tail and abbreviate \(x_0=x_r\), \(A=A_{\rho,K}\).  Put
\begin{equation}
  n=\lfloor\mu-x_0\rfloor,
  \qquad q=x_0+n,
  \qquad B=\lceil K-x_0\rceil.
  \label{eq:sc-interface-bookkeeping}
\end{equation}
Then
\[
  q\le\mu<q+1,
  \qquad n<B,
  \qquad G_K(x_0+B)=0.
\]
The point \(q\) is the largest half-integer in the tail grid that does not
exceed \(\mu\); it need not equal \(\lfloor\mu\rfloor\).  A low tail also
forces \(\mu>x_0\ge1/2\), hence \(\mu>1/2\).

Let
\[
  h_A(t)=A(x_0+t)+\frac14,
  \qquad h_K(t)=G_K(x_0+t)+\frac14.
\]
All samples at and beyond \(B\) have floor zero, so
\[
  \frac{\mathcal T_r}{2}=\mathbf T(h_A;0,B).
\]
If \(n\ge1\), additivity of the trapezoidal sum at the shared endpoint and
\(A\le G_K\) give
\begin{equation}
  \frac{\mathcal T_r}{2}
  \le\mathbf T(h_A;0,n)+\mathbf T(h_K;n,B).
  \label{eq:sc-tail-split-positive-n}
\end{equation}
At the split point, the two half-weights add to the original full sample;
after replacing the suffix by \(G_K\), its half-sample can only increase.
If \(n=0\), there is no concave block, and instead
\begin{equation}
  \frac{\mathcal T_r}{2}\le\mathbf T(h_K;0,B).
  \label{eq:sc-tail-split-zero-n}
\end{equation}

For \(0<z<\mu\),
\begin{equation}
  A''(z)=\frac1{\pi\sqrt{K^2-z^2}}
  -\frac1{\pi\sqrt{\mu^2-z^2}}<0.
  \label{eq:sc-action-concavity}
\end{equation}
Thus \(A\) is concave on \([x_0,q]\), also when \(q=\mu\) by continuity.
For \(n\ge1\), the concave bound
\eqref{eq:sc-concave-floor-bound} gives
\begin{equation}
  \mathbf T(h_A;0,n)
  \le\int_0^nA(x_0+t)\,dt+\frac n4.
  \label{eq:sc-concave-head}
\end{equation}

For the suffix set \(g(t)=G_K(x_0+t)\), using the zero extension after
\(K\).  It is nonnegative, decreasing, convex, globally
\(1/2\)-Lipschitz, and \(g(B)=0\).  Because
\(q\le\mu=\rho K<\omega_0K<K/2\),
\[
  t_*:=\frac K2-x_0
\]
belongs to \([n,B]\) (strictly to the right of \(n\)).  On
\([t_*,B]\),
\[
  |g'(t)|\le\frac1\pi\arccos\frac12=\frac13,
\]
and
\[
  g(t_*)=G_K(K/2)=\omega_0K.
\]
Every hypothesis of \eqref{eq:sc-improved-convex-floor-bound} is therefore
satisfied, and
\begin{equation}
  \mathbf T(h_K;n,B)
  \le\int_n^BG_K(x_0+t)\,dt
  -\frac14\lfloor\omega_0K\rfloor.
  \label{eq:sc-convex-suffix}
\end{equation}
The identical statement with \(n=0\) applies to
\eqref{eq:sc-tail-split-zero-n}.

Combining \eqref{eq:sc-concave-head} and
\eqref{eq:sc-convex-suffix}, or using only the latter when \(n=0\), leaves
one and only one integral mismatch:
\begin{equation}
  \delta:=\int_q^\mu G_\mu(z)\,dz.
  \label{eq:sc-interface-loss}
\end{equation}
Indeed,
\[
  \int_{x_0}^{q}A(z)\,dz+\int_q^KG_K(z)\,dz
  =\int_{x_0}^KA(z)\,dz+\delta,
\]
and the same identity holds for \(n=0\), where \(q=x_0\).  Hence
\begin{equation}
  \frac{\mathcal T_r}{2}
  \le\int_{x_0}^KA(z)\,dz
  +\delta-\frac{\lfloor\omega_0K\rfloor-n}{4}.
  \label{eq:sc-tail-compensation}
\end{equation}

The turning-point estimate of Lemma~6.2 in \cite{FLPSannuli} states that
\begin{equation}
  G_\mu(z)<\frac{(\mu-z)^{3/2}}{3\sqrt\mu},
  \qquad 0<z<\mu.
  \label{eq:sc-turning-point-bound}
\end{equation}
If \(q<\mu\), integration gives a strict first bound; if \(q=\mu\), both
sides below vanish.  The endpoint-safe statement is
\begin{equation}
  0\le\delta
  \le\frac{2(\mu-q)^{5/2}}{15\sqrt\mu}
  <\frac{2\sqrt2}{15},
  \label{eq:sc-interface-loss-bound}
\end{equation}
because \(0\le\mu-q<1\) and \(\mu>1/2\).  The last inequality remains
strict when \(q=\mu\), since then \(\delta=0\).

Finally, \(x_0\ge1/2\) implies
\[
  n\le\mu-x_0\le\rho K-\frac12,
\]
while \(\lfloor y\rfloor>y-1\) for every real \(y\).  Therefore
\begin{align*}
  \lfloor\omega_0K\rfloor-n
  &>\omega_0K-1-\left(\rho K-\frac12\right)\\
  &=K(\omega_0-\rho)-\frac12\\
  &\ge\frac{8\sqrt2}{15}
  >4\delta.
\end{align*}
Both strict inequalities remain strict when the threshold in
\eqref{eq:sc-low-interface-sector} is attained.  Substitution in
\eqref{eq:sc-tail-compensation} proves
\eqref{eq:sc-low-tail-conclusion}.
\end{proof}

\subsection{Uniform closure of the small-hole endpoint}

Set
\begin{equation}
  \rho_*:=\frac{\omega_0}{1+2C_*}
  =\frac{\frac{\sqrt3}{2\pi}-\frac16}
  {2+\frac{16\sqrt2}{15}}.
  \label{eq:sc-rho-star}
\end{equation}
Since \(\omega_0>0\) and \(C_*>0\),
\(0<\rho_*<\omega_0\).  Its defining identity is
\begin{equation}
  \omega_0-\rho_*=2C_*\rho_*,
  \qquad
  \frac{\omega_0-\rho_*}{2\rho_*}=C_*.
  \label{eq:sc-rho-star-identity}
\end{equation}

\begin{theorem}[Complete small-hole endpoint]
\label{thm:sc-small-hole-endpoint}
For every \(0<\rho\le\rho_*\) and every \(K\ge0\),
\begin{equation}
  \boxed{
  N_D(A_{\rho,1},K^2)
  \le\frac{2}{9\pi}(1-\rho^3)K^3.}
  \label{eq:sc-small-hole-polya}
\end{equation}
The face \(\rho=\rho_*\), the interface wall \(\rho K=1/2\), and every
strict spectral wall are included.
\end{theorem}

\begin{proof}
If \(K=0\), the strict spectrum is positive, so the left side and the right
side of \eqref{eq:sc-small-hole-polya} are both zero.  Assume \(K>0\) and
write \(\mu=\rho K\).

First suppose \(\mu\le1/2\).  Every shifted-tail start satisfies
\[
  x_r=r+\frac12\ge\frac12\ge\mu.
\]
Lemma~\ref{lem:sc-high-interface-tails} therefore proves
\eqref{eq:sc-tail-target} for every active tail.  At the joining wall
\(\mu=1/2\), the first start is exactly \(x_0=\mu\), which belongs to the
inclusive high-interface lemma.

Now suppose \(\mu>1/2\).  Then
\begin{equation}
  K>\frac1{2\rho}.
  \label{eq:sc-high-K-from-interface}
\end{equation}
For \(0<\rho\le\rho_*\), the defining identity
\eqref{eq:sc-rho-star-identity} gives
\begin{equation}
  \frac{\omega_0-\rho}{2\rho}\ge C_*;
  \label{eq:sc-uniform-small-hole-threshold}
\end{equation}
equivalently, multiply
\((1+2C_*)\rho\le(1+2C_*)\rho_*=\omega_0\)
by positive quantities.  Since
\(\omega_0-\rho>0\), \eqref{eq:sc-high-K-from-interface} and
\eqref{eq:sc-uniform-small-hole-threshold} imply
\begin{equation}
  K(\omega_0-\rho)
  >\frac{\omega_0-\rho}{2\rho}
  \ge C_*.
  \label{eq:sc-strict-small-hole-threshold}
\end{equation}
Every start \(x_r\ge\mu\) is controlled by
Lemma~\ref{lem:sc-high-interface-tails}; every start \(x_r<\mu\) is
controlled by Lemma~\ref{lem:sc-low-interface-tails}, because
\(\rho\le\rho_*<\omega_0\) and
\eqref{eq:sc-strict-small-hole-threshold} supplies its threshold.  Thus all
active tails satisfy \eqref{eq:sc-tail-target} in this branch as well.

Lemma~\ref{lem:sc-weighted-scaffold} now yields
\[
  S_{\rho,K}\le\frac{2}{9\pi}(1-\rho^3)K^3,
\]
and Proposition~\ref{prop:sc-phase-lattice} gives
\(N_D(A_{\rho,1},K^2)\le S_{\rho,K}\), proving
\eqref{eq:sc-small-hole-polya}.

At \(\rho=\rho_*\), inequality
\eqref{eq:sc-uniform-small-hole-threshold} is an equality, but the first
inequality in \eqref{eq:sc-strict-small-hole-threshold} is strict whenever
\(\mu>1/2\).  Hence the endpoint is included.  If instead
\(K(\omega_0-\rho)=C_*\), then
\[
  \mu=\frac{C_*\rho}{\omega_0-\rho}\le\frac12
\]
by \eqref{eq:sc-uniform-small-hole-threshold}, so the point already belongs
to the all-high-interface branch.  No limiting argument at \(\rho=0\), no
continuity of Bessel roots, and no finite numerical verification is used.
\end{proof}

\section{The thin-shell endpoint}
\label{thin:sec-endpoint}

We now prove the P\'olya inequality uniformly as the shell becomes thin.  In
this section the counting function is strict:
\[
  N_D(\Omega,K^2)=\#\{j:\lambda_j(\Omega)<K^2\}.
\]
We retain the global convention
\([u]_<:=\#\{n\in\mathbb N:n<u\}=\max\{0,\lceil u\rceil-1\}\).
Thus a spectral level lying exactly on the wall \(K^2\) is not counted.  Put
\begin{equation}
  \varepsilon=1-\rho,
  \qquad a=\varepsilon K,
  \qquad
  m_\varepsilon=1-\varepsilon+\frac{\varepsilon^2}{3}.
  \label{thin:eq-normalized-variables}
\end{equation}
Since \(1-\rho^3=3\varepsilon m_\varepsilon\), the desired right-hand side
has the two equivalent forms
\begin{equation}
  \frac{2}{9\pi}(1-\rho^3)K^3
  =\frac{1}{\varepsilon^2}\frac{2a^3}{3\pi}m_\varepsilon.
  \label{thin:eq-weyl-normalization}
\end{equation}

The proof uses the optimized phase estimate already established in
Proposition~\ref{prop:sc-phase-lattice}.  We record precisely the consequence
that is used here.  For
\(\nu_\ell=\ell+\tfrac12\), let
\(\gamma_{\rho,K}(\nu_\ell)\) be the normalized radial Bessel phase.  If
\(0\leq\nu_\ell\leq K\), that proposition gives
\begin{equation}
  \gamma_{\rho,K}(\nu_\ell)
  < A_{\rho,K}(\nu_\ell)+\frac14,
  \label{thin:eq-phase-input}
\end{equation}
where, after the scaling in \eqref{thin:eq-normalized-variables},
\(A_{\rho,K}(\nu)=\mathcal A(\varepsilon\nu)\) and
\(\mathcal A\) is the explicit action in
\eqref{thin:eq-action-definition} below.  The exact radial count in channel
\(\ell\) is
\begin{equation}
  \#\{n\geq1:n<\gamma_{\rho,K}(\nu_\ell)\}.
  \label{thin:eq-exact-phase-count}
\end{equation}
In particular, \eqref{thin:eq-phase-input} is an inequality for the strict
count; it is not a replacement of that count by an ordinary floor.

\subsection{The product comparison}

Under the standard unitary radial transformation, the angular channel
\(\ell\in\mathbb N_0\) is the Dirichlet operator
\[
  L_\ell=-\frac{d^2}{dr^2}+\frac{\ell(\ell+1)}{r^2}
  \quad\hbox{in }L^2((\rho,1),dr),
\]
and it occurs with multiplicity \(2\ell+1\).  The following lemma includes
the direction of the min--max comparison and every strict radial and angular
wall.

\begin{lemma}[Strict product majorant]
\label{thin:lem-product-majorant}
Let \(0<\varepsilon<1\), \(a=\varepsilon K\), and define
\begin{equation}
  N(a)=\max\left\{0,\left\lceil\frac a\pi\right\rceil-1\right\}.
  \label{thin:eq-radial-count}
\end{equation}
For \(1\leq n\leq N(a)\), put
\[
  b_n=\sqrt{a^2-n^2\pi^2},
  \qquad
  M_n=
  \left\lceil
  \sqrt{\left(\frac{b_n}{\varepsilon}\right)^2+\frac14}
  -\frac12
  \right\rceil.
\]
Then
\begin{equation}
  N_D(A_{\rho,1},K^2)
  \leq \mathcal P_\varepsilon(a)
  :=\sum_{n=1}^{N(a)}M_n^2.
  \label{thin:eq-product-majorant}
\end{equation}
In particular, the strict shell count is zero for \(0\leq a\leq\pi\).
\end{lemma}

\begin{proof}
The two quadratic forms have the common domain \(H_0^1(\rho,1)\), and
\[
  \int_\rho^1\left(
     |u'|^2+\frac{\ell(\ell+1)}{r^2}|u|^2
  \right)dr
  \geq
  \int_\rho^1\left(|u'|^2+\ell(\ell+1)|u|^2\right)dr,
\]
because \(r\leq1\).  The min--max principle therefore gives
\[
  \lambda_{\ell,n}
  \geq \left(\frac{n\pi}{\varepsilon}\right)^2
       +\ell(\ell+1),
  \qquad n\geq1.
\]
Consequently, a shell eigenvalue in this channel can be counted below
\(K^2\) only if
\begin{equation}
  n^2\pi^2+\varepsilon^2\ell(\ell+1)<a^2.
  \label{thin:eq-product-condition}
\end{equation}
The active radial integers are exactly those satisfying \(n\pi<a\), whose
number is \eqref{thin:eq-radial-count}.  For a fixed active \(n\), condition
\eqref{thin:eq-product-condition} is
\[
  \ell<
  \sqrt{\left(\frac{b_n}{\varepsilon}\right)^2+\frac14}-\frac12.
\]
Hence the allowed angular indices are
\(0,\ldots,M_n-1\), and their full spherical multiplicity is
\[
  \sum_{\ell=0}^{M_n-1}(2\ell+1)=M_n^2.
\]
Summing proves \eqref{thin:eq-product-majorant}.

If \(a=m\pi\), the radial level \(n=m\) is excluded by the strict
inequality \(n\pi<a\); thus \(N(a)=m-1\).  At an angular wall
\(b_n^2/\varepsilon^2=L(L+1)\), the displayed expression inside the ceiling
is exactly \(L\), so \(M_n=L\) and the level \(\ell=L\) is excluded.  The
jump of multiplicity \(2L+1\) occurs only above the wall.  Finally, no
positive integer satisfies \(n\pi<a\) when \(0\leq a\leq\pi\), including
both endpoints, so the product majorant and hence the shell count are zero
there.
\end{proof}

\begin{lemma}[Exact product deficit]
\label{thin:lem-product-deficit}
For \(a>\pi\), let \(N=N(a)\) and write uniquely
\begin{equation}
  \frac a\pi=N+t,
  \qquad N\geq1,
  \qquad 0<t\leq1.
  \label{thin:eq-strict-radial-parametrization}
\end{equation}
At \(a=m\pi\), this means \((N,t)=(m-1,1)\).  Define
\[
  I_0(a)=\frac{2a^3}{3\pi},
  \qquad
  S_0(a)=\sum_{n=1}^{N}(a^2-n^2\pi^2),
  \qquad
  D(a)=I_0(a)-S_0(a).
\]
Then
\begin{equation}
  \frac{D(a)}{\pi^2}
  =\frac{N^2}{2}+N\left(t^2+\frac16\right)+\frac{2t^3}{3},
  \label{thin:eq-deficit-identity}
\end{equation}
and, uniformly for \(a>\pi\),
\begin{equation}
  D(a)>\frac25a^2.
  \label{thin:eq-deficit-bound}
\end{equation}
\end{lemma}

\begin{proof}
The sum of the first \(N\) squares gives
\[
\begin{aligned}
  \frac{D(a)}{\pi^2}
  &=\frac23(N+t)^3-N(N+t)^2
    +\frac{N(N+1)(2N+1)}6 \\
  &=\frac{N^2}{2}+N\left(t^2+\frac16\right)+\frac{2t^3}{3},
\end{aligned}
\]
which proves \eqref{thin:eq-deficit-identity}.  At \(a=m\pi\), its value
with \((N,t)=(m-1,1)\) agrees with the right limit \((N,t)\to(m,0+)\);
both equal \(m^2/2+m/6\).  Thus the formula owns the wall and both
one-sided limits.

Subtracting \(\tfrac25(N+t)^2\) from
\eqref{thin:eq-deficit-identity} gives the exact identity
\begin{equation}
\begin{aligned}
  \frac{D(a)}{\pi^2}-\frac25\left(\frac a\pi\right)^2
  ={}&\frac{N^2}{10}
  +N\left[\left(t-\frac25\right)^2+\frac1{150}\right] \\
  &+\frac{2t^3}{3}-\frac{2t^2}{5}.
\end{aligned}
  \label{thin:eq-deficit-square-completion}
\end{equation}
For \(h(t)=2t^3/3-2t^2/5\),
\(h'(t)=2t(t-2/5)\), so the minimum on \(0<t\leq1\) is
\(h(2/5)=-8/375\).  Therefore the right-hand side of
\eqref{thin:eq-deficit-square-completion} is at least
\[
  \frac{N^2}{10}+\frac{N}{150}-\frac8{375}.
\]
This expression increases for integral \(N\geq1\), and at \(N=1\) it is
\[
  \frac1{10}+\frac1{150}-\frac8{375}
  =\frac{32}{375}>0.
\]
This proves \eqref{thin:eq-deficit-bound} without sampling.
\end{proof}

\begin{proposition}[The inclusive low-action piece]
\label{thin:prop-low-piece}
If
\[
  0<\varepsilon\leq\frac18,
  \qquad 0\leq a\leq\frac{\pi}{4\varepsilon},
\]
then
\begin{equation}
  N_D(A_{\rho,1},K^2)
  <\frac{1}{\varepsilon^2}\frac{2a^3}{3\pi}m_\varepsilon
  \qquad(a>0),
  \label{thin:eq-low-piece-conclusion}
\end{equation}
whereas both sides are zero at \(a=0\).
\end{proposition}

\begin{proof}
For \(0<a\leq\pi\), Lemma~\ref{thin:lem-product-majorant} gives a zero
left-hand side and the right-hand side of
\eqref{thin:eq-low-piece-conclusion} is positive.  It remains to consider
\(a>\pi\).

For \(B=b_n/\varepsilon>0\), the elementary strict ceiling inequality gives
\[
  M_n<\sqrt{B^2+\frac14}+\frac12.
\]
Both sides are nonnegative, and
\(\sqrt{B^2+1/4}\leq B+1/2\).  Squaring and restoring the scale yields
\begin{equation}
  \varepsilon^2M_n^2
  <b_n^2+\varepsilon b_n+\varepsilon^2.
  \label{thin:eq-ceiling-cost}
\end{equation}
The function \(x\mapsto\sqrt{a^2-\pi^2x^2}\) is strictly decreasing on
its positive interval.  Hence
\begin{equation}
  \sum_{n=1}^{N}b_n
  <\int_0^N\sqrt{a^2-\pi^2x^2}\,dx
  \leq\int_0^{a/\pi}\sqrt{a^2-\pi^2x^2}\,dx
  =\frac{a^2}{4}.
  \label{thin:eq-quarter-disc}
\end{equation}
Also \(N<a/\pi\), including every radial wall because the wall level is
excluded.  Summation of \eqref{thin:eq-ceiling-cost} therefore gives
\begin{equation}
  \varepsilon^2\mathcal P_\varepsilon(a)
  <S_0(a)+\frac{\varepsilon a^2}{4}
    +\frac{\varepsilon^2a}{\pi}.
  \label{thin:eq-summed-ceiling-cost}
\end{equation}

It is enough to prove
\begin{equation}
  D(a)\geq
  \varepsilon\left(I_0(a)+\frac{a^2}{4}\right)
  +\frac{\varepsilon^2a}{\pi},
  \label{thin:eq-sufficient-deficit}
\end{equation}
because then \eqref{thin:eq-summed-ceiling-cost} is at most
\(I_0(a)(1-\varepsilon)\), and hence is strictly less than
\(I_0(a)m_\varepsilon\).

After division by \(a^2\), the right-hand side of
\eqref{thin:eq-sufficient-deficit} is
\[
  \frac{2\varepsilon a}{3\pi}+\frac\varepsilon4
  +\frac{\varepsilon^2}{\pi a}.
\]
The low-piece inequalities \(a\leq\pi/(4\varepsilon)\) and \(a\geq\pi\),
together with the elementary strict bound \(\pi>3\), imply
\begin{equation}
  \frac{2\varepsilon a}{3\pi}+\frac\varepsilon4
  +\frac{\varepsilon^2}{\pi a}
  \leq \frac16+\frac\varepsilon4+\frac{\varepsilon^2}{9}.
  \label{thin:eq-low-cost-screen}
\end{equation}
The last expression increases with \(\varepsilon>0\).  On the closed face
\(\varepsilon=1/8\), its exact reserve below \(2/5\) is
\begin{equation}
  \frac25-\left(\frac16+\frac1{32}+\frac1{576}\right)
  =\frac{577}{2880}>0.
  \label{thin:eq-low-reserve}
\end{equation}
Thus Lemma~\ref{thin:lem-product-deficit} proves
\eqref{thin:eq-sufficient-deficit}, and
Lemma~\ref{thin:lem-product-majorant} together with
\eqref{thin:eq-summed-ceiling-cost} proves
\eqref{thin:eq-low-piece-conclusion}.

All inequalities allow \(a=\pi/(4\varepsilon)\).  Since
\(\varepsilon\leq1/8\), this face has \(a\geq2\pi\).  At its corner
\((\varepsilon,a)=(1/8,2\pi)\), the strict radial convention is
\((N,t)=(1,1)\), and the positive reserve
\eqref{thin:eq-low-reserve} remains unchanged.  Thus the low proof owns the
whole common face.
\end{proof}

\subsection{The complementary-action comparison}

For \(r\in\{\rho,1\}\) and \(0\leq y\leq ar\), define
\begin{equation}
  J_r(y)=\sqrt{a^2r^2-y^2}
          -y\arccos\frac{y}{ar},
  \label{thin:eq-J-action}
\end{equation}
and set
\begin{equation}
  \mathcal A(y)=
  \begin{cases}
  \dfrac{J_1(y)-J_\rho(y)}{\pi\varepsilon},
    &0\leq y\leq\rho a,\\[6pt]
  \dfrac{J_1(y)}{\pi\varepsilon},
    &\rho a\leq y\leq a.
  \end{cases}
  \label{thin:eq-action-definition}
\end{equation}
We extend \(\mathcal A(y)=0\) for \(y>a\).

\begin{lemma}[Action geometry and curvature switch]
\label{thin:lem-action-geometry}
The function \(\mathcal A\) is continuous, continuously differentiable at
\(y=\rho a\), and strictly decreasing from
\[
  T:=\mathcal A(0)=\frac a\pi
  \quad\hbox{to}\quad \mathcal A(a)=0.
\]
Let \(R:[0,T]\to[0,a]\) be its decreasing inverse, let
\[
  \tau=\mathcal A(\rho a),
  \qquad F(t)=R(t)^2,
  \qquad g(t)=-F'(t).
\]
Then
\begin{equation}
\begin{gathered}
  g\text{ decreases on }(0,\tau),
  \qquad g\text{ increases on }(\tau,T),\\
  F(0)=a^2,
  \qquad F(\tau)=\rho^2a^2,
  \qquad F(T)=0,
  \qquad g(T)=2\pi\rho a,
\end{gathered}
  \label{thin:eq-action-geometry}
\end{equation}
where \(g(T)\) denotes the one-sided limit.  Moreover \(g(t)=O(t^{-1/3})\)
as \(t\downarrow0\).
\end{lemma}

\begin{proof}
Direct differentiation on the appropriate open intervals gives
\begin{equation}
  J_r'(y)=-\arccos\frac{y}{ar},
  \qquad
  J_r''(y)=\frac1{\sqrt{a^2r^2-y^2}}.
  \label{thin:eq-J-derivatives}
\end{equation}
At \(y=\rho a\), \(J_\rho(\rho a)=0\), and the two one-sided first
derivatives of \eqref{thin:eq-action-definition} both equal
\(-\arccos\rho/(\pi\varepsilon)\).  Thus the two pieces and their first
derivatives match.  Both derivatives are negative away from \(y=a\), so the
inverse exists on exactly the stated interval.

On the outer branch \(\rho a<y<a\), put
\(\alpha=\arccos(y/a)\).  If \(t=\mathcal A(y)\), then
\begin{equation}
  g(t)=-\frac{2y}{\mathcal A'(y)}
      =\frac{2\pi\varepsilon y}{\alpha}.
  \label{thin:eq-g-outer}
\end{equation}
The function \(y/\alpha\) increases with \(y\), whereas \(y=R(t)\)
decreases with \(t\); hence \(g\) decreases on \(0<t<\tau\).

On the inner branch \(0<y<\rho a\), put
\[
  \delta(y)=\arccos\frac ya-\arccos\frac{y}{\rho a}>0.
\]
Differentiating \(\arccos(y/(ar))\) with respect to \(r\) gives the exact
identity
\begin{equation}
  \frac{\delta(y)}{y}
  =\int_\rho^1\frac{dr}{r\sqrt{a^2r^2-y^2}}.
  \label{thin:eq-delta-integral}
\end{equation}
The integrand increases strictly with \(y\), so \(\delta(y)/y\) increases,
and \(y/\delta(y)\) decreases, with \(y\).  Since
\[
  g(t)=\frac{2\pi\varepsilon y}{\delta(y)}
\]
and \(y=R(t)\) decreases, \(g\) increases on \(\tau<t<T\).  Formula
\eqref{thin:eq-g-outer} and its inner counterpart agree at the interface.

The three values of \(F\) follow from
\(R(0)=a\), \(R(\tau)=\rho a\), and \(R(T)=0\).  As \(y\downarrow0\),
\eqref{thin:eq-delta-integral} gives
\[
  \delta(y)=\frac{\varepsilon}{\rho a}y+O(y^3),
\]
and therefore \(g(T)=2\pi\rho a\).

Finally, on the outer branch write \(y=a\cos\theta\).  Then
\begin{equation}
  t=\frac{a}{\pi\varepsilon}
       (\sin\theta-\theta\cos\theta)
    \sim\frac{a\theta^3}{3\pi\varepsilon},
  \qquad
  g(t)=\frac{2\pi\varepsilon a\cos\theta}{\theta}
       =O(t^{-1/3}).
  \label{thin:eq-improper-asymptotic}
\end{equation}
This proves the asserted improper endpoint behavior.  All square roots and
arccosines in the proof are taken on their stated closed domains; every
division occurs in an open interval where its denominator is positive, and
the endpoint zeros are handled by the limits just computed.
\end{proof}

\begin{lemma}[Strict phase-to-action layer cake]
\label{thin:lem-layer-cake}
For \(\ell\in\mathbb N_0\), set
\begin{equation}
  y_\ell=\varepsilon\left(\ell+\frac12\right),
  \qquad
  q_\ell=\left[\mathcal A(y_\ell)+\frac14\right]_<,
  \qquad
  P_{\mathcal A}=\varepsilon\sum_{\ell\geq0}y_\ell q_\ell.
  \label{thin:eq-phase-proxy}
\end{equation}
Then
\begin{equation}
  N_D(A_{\rho,1},K^2)
  \leq\frac{2P_{\mathcal A}}{\varepsilon^2}.
  \label{thin:eq-phase-transfer}
\end{equation}
For \(n\geq1\), put \(x_n=n-\tfrac14\), and define
\begin{equation}
  M_\varepsilon(x)
  =\#\left\{\ell\in\mathbb N_0:
       \varepsilon\left(\ell+\frac12\right)<x\right\}
  =\max\left\{0,
       \left\lceil\frac{x}{\varepsilon}-\frac12\right\rceil
       \right\}.
  \label{thin:eq-half-integer-count}
\end{equation}
The proxy has the exact finite representation
\begin{equation}
  P_{\mathcal A}
  =\frac{\varepsilon^2}{2}
    \sum_{\substack{n\geq1\\x_n<T}}
       M_\varepsilon(R(x_n))^2.
  \label{thin:eq-discrete-layer-cake}
\end{equation}
Moreover,
\begin{equation}
  I:=\frac12\int_0^T F(t)\,dt
  =\frac{a^3}{3\pi}
     \left(1-\varepsilon+\frac{\varepsilon^2}{3}\right),
  \qquad
  \frac{2I}{\varepsilon^2}
  =\frac{2}{9\pi}(1-\rho^3)K^3.
  \label{thin:eq-continuous-normalization}
\end{equation}
\end{lemma}

\begin{proof}
For \(y_\ell<a\), equations \eqref{thin:eq-phase-input} and
\eqref{thin:eq-exact-phase-count} show that the number of radial levels in
channel \(\ell\) is at most \(q_\ell\).  If \(y_\ell\geq a\), equivalently
\(\nu_\ell\geq K\), the product comparison gives
\[
  \lambda_{\ell,1}
  \geq\frac{\pi^2}{\varepsilon^2}+\ell(\ell+1)
  =\frac{\pi^2}{\varepsilon^2}+\nu_\ell^2-\frac14
  >K^2,
\]
so there is no radial level; at \(y_\ell=a\), the definition also gives
\(q_\ell=[1/4]_< =0\).  Multiplicity \(2\ell+1=2y_\ell/\varepsilon\)
now yields
\[
  N_D(A_{\rho,1},K^2)
  \leq\sum_{\ell\geq0}\frac{2y_\ell}{\varepsilon}q_\ell
  =\frac{2P_{\mathcal A}}{\varepsilon^2},
\]
which proves \eqref{thin:eq-phase-transfer}.  If
\(\mathcal A(y_\ell)+1/4=m\in\mathbb Z\), then \(q_\ell=m-1\), so the
phase-proxy wall is excluded exactly as required.

For every \(y\geq0\), the strict bracket is
\[
  \left[\mathcal A(y)+\frac14\right]_<
  =\#\{n\geq1:x_n<\mathcal A(y)\}.
\]
The sums are finite.  Since \(\mathcal A\) and \(R\) are decreasing,
\(x_n<\mathcal A(y_\ell)\) is equivalent to \(y_\ell<R(x_n)\).  Therefore
\[
\begin{aligned}
  P_{\mathcal A}
  &=\varepsilon
    \sum_{\substack{n\geq1\\x_n<T}}
      \sum_{y_\ell<R(x_n)}y_\ell \\
  &=\varepsilon
    \sum_{\substack{n\geq1\\x_n<T}}
      \varepsilon\sum_{\ell=0}^{M_\varepsilon(R(x_n))-1}
        \left(\ell+\frac12\right) \\
  &=\frac{\varepsilon^2}{2}
    \sum_{\substack{n\geq1\\x_n<T}}
      M_\varepsilon(R(x_n))^2,
\end{aligned}
\]
because \(\sum_{\ell=0}^{M-1}(\ell+1/2)=M^2/2\).  This proves
\eqref{thin:eq-discrete-layer-cake}.  Formula
\eqref{thin:eq-half-integer-count} follows from the same strict counting;
at \(x/\varepsilon=m+1/2\), it equals \(m\), excluding \(\ell=m\).

Differentiation of \eqref{thin:eq-J-action} with respect to \(r\) gives
\[
  \frac{\partial J_r(y)}{\partial r}
  =\sqrt{a^2-\frac{y^2}{r^2}}
  \quad(0\leq y\leq ar).
\]
Consequently the two cases of \eqref{thin:eq-action-definition} combine as
\begin{equation}
  \mathcal A(y)
  =\frac1{\pi\varepsilon}
    \int_\rho^1
      \left(a^2-\frac{y^2}{r^2}\right)_+^{1/2}\,dr,
  \label{thin:eq-action-integral}
\end{equation}
where \(s_+=\max\{s,0\}\).
The area identity for a decreasing inverse and Tonelli's theorem give
\[
  \frac12\int_0^T R(t)^2\,dt
  =\int_0^a y\mathcal A(y)\,dy.
\]
Using \eqref{thin:eq-action-integral} and then \(y=rz\),
\[
\begin{aligned}
  \int_0^a y\mathcal A(y)\,dy
  &=\frac1{\pi\varepsilon}
    \int_\rho^1r^2\,dr
    \int_0^a z\sqrt{a^2-z^2}\,dz \\
  &=\frac{a^3}{3\pi\varepsilon}\int_\rho^1r^2\,dr
   =\frac{a^3}{3\pi}
      \left(1-\varepsilon+\frac{\varepsilon^2}{3}\right).
\end{aligned}
\]
This is \eqref{thin:eq-continuous-normalization}; its final equality follows
from \(1-\rho^3=3\varepsilon m_\varepsilon\) and \(K=a/\varepsilon\).
\end{proof}

\begin{lemma}[Shifted sawtooth and radial quadrature deficit]
\label{thin:lem-radial-deficit}
Let
\[
  C(t)=\#\{n\geq1:n-\tfrac14<t\},
  \qquad w(t)=t-C(t),
  \qquad \mathfrak W(t)=\int_0^t w(s)\,ds,
\]
and put \(\Psi(t)=\mathfrak W(t)-t/4\).  Then, for all \(t\geq0\),
\begin{equation}
  \mathfrak W(t)\geq0,
  \qquad
  -\frac1{32}\leq\Psi(t)\leq\frac3{32}.
  \label{thin:eq-sawtooth-bounds}
\end{equation}
Furthermore, with
\begin{equation}
  D_{\mathrm{rad}}
  =\int_0^T F(t)\,dt
   -\sum_{\substack{n\geq1\\x_n<T}}F(x_n),
  \label{thin:eq-radial-deficit-definition}
\end{equation}
one has
\begin{equation}
  D_{\mathrm{rad}}
  \geq\frac{\rho^2a^2}{4}-\frac{\pi\rho a}{4}.
  \label{thin:eq-radial-deficit-bound}
\end{equation}
\end{lemma}

\begin{proof}
On \(0\leq t\leq3/4\), \(C(t)=0\), and hence
\begin{equation}
  \Psi(t)=\frac12\left(t-\frac14\right)^2-\frac1{32}.
  \label{thin:eq-first-sawtooth-cell}
\end{equation}
For \(x_n=n-1/4\), \(n\geq1\), and \(t=x_n+u\) with
\(0\leq u\leq1\), direct integration gives
\begin{equation}
  \Psi(t)=\frac12\left(u-\frac12\right)^2-\frac1{32}.
  \label{thin:eq-periodic-sawtooth-cell}
\end{equation}
These formulas agree at cell endpoints and prove the two-sided bound for
\(\Psi\).  On the first cell \(\mathfrak W=t^2/2\geq0\).  Thereafter
\(\mathfrak W=t/4+\Psi\geq t/4-1/32>0\), proving
\eqref{thin:eq-sawtooth-bounds}.  At a jump \(x_n\), strict counting gives
the one-sided values \(w(x_n^-)=3/4\) and \(w(x_n^+)=-1/4\), while
\(\mathfrak W\)
and \(\Psi\) are continuous.

Distributionally, \(dw=dt-dC\).  Since \(F(T)=0\), an atom at \(x_n=T\)
has zero weight and is consistent with its exclusion from
\eqref{thin:eq-radial-deficit-definition}.  Stieltjes integration by parts,
first on \([\eta,T]\) and then with \(\eta\downarrow0\), gives
\begin{equation}
  D_{\mathrm{rad}}=\int_0^T w(t)g(t)\,dt.
  \label{thin:eq-stieltjes-start}
\end{equation}
Indeed, \eqref{thin:eq-improper-asymptotic}, together with
\(w(t)=t\) and \(\mathfrak W(t)=t^2/2\) near zero, gives
\(w(t)g(t)=O(t^{2/3})\) and
\(\mathfrak W(t)g(t)=O(t^{5/3})\to0\), so the improper
boundary term vanishes.

Split \eqref{thin:eq-stieltjes-start} at the actual value \(\tau\), with
no assumption that \(\tau\) belongs to any particular sawtooth cell.  On
\((0,\tau)\), \(g\) decreases and \(\mathfrak W\geq0\); hence
\begin{equation}
  \int_0^\tau wg
  =\mathfrak W(\tau)g(\tau)-\int_0^\tau \mathfrak W\,dg
  \geq \mathfrak W(\tau)g(\tau).
  \label{thin:eq-decreasing-stieltjes}
\end{equation}
On \((\tau,T)\), \(g\) increases and \(w=1/4+\Psi'\) almost everywhere.
Adding its integration-by-parts identity to
\eqref{thin:eq-decreasing-stieltjes}, and using
\(\mathfrak W(\tau)=\tau/4+\Psi(\tau)\), cancels the interface value and
yields
\[
  D_{\mathrm{rad}}
  \geq \frac{F(\tau)}4+\frac{\tau g(\tau)}4
     +\Psi(T)g(T)-\int_\tau^T\Psi\,dg.
\]
Because \(dg\) is a positive measure on \((\tau,T)\), the bounds
\eqref{thin:eq-sawtooth-bounds} imply
\[
\begin{aligned}
  D_{\mathrm{rad}}
  &\geq \frac{F(\tau)}4-\frac{g(T)}8
    +g(\tau)\left(\frac\tau4+\frac3{32}\right) \\
  &\geq \frac{F(\tau)}4-\frac{g(T)}8.
\end{aligned}
\]
The two discarded quantities,
\(-\int_0^\tau \mathfrak W\,dg\) and
\(g(\tau)(\tau/4+3/32)\), are explicitly nonnegative.  Substitution of
\(F(\tau)=\rho^2a^2\) and \(g(T)=2\pi\rho a\) from
Lemma~\ref{thin:lem-action-geometry} proves
\eqref{thin:eq-radial-deficit-bound}.  Continuity of
\(\mathfrak W,\Psi,g\) at the
split shows that the same proof covers \(\tau\) below, on, or above every
sawtooth grid point.
\end{proof}

\begin{proposition}[The inclusive high-action piece]
\label{thin:prop-high-piece}
If
\[
  0<\varepsilon\leq\frac18,
  \qquad a\geq\frac{\pi}{4\varepsilon},
\]
then
\begin{equation}
  N_D(A_{\rho,1},K^2)
  <\frac{2}{9\pi}(1-\rho^3)K^3.
  \label{thin:eq-high-piece-conclusion}
\end{equation}
\end{proposition}

\begin{proof}
The number of active radial action layers is
\begin{equation}
\begin{aligned}
  J_{\mathrm{rad}}
  &:=\#\{n\geq1:x_n<T\}
    =\left\lceil T+\frac14\right\rceil-1 \\
  &<T+\frac14=\frac a\pi+\frac14.
\end{aligned}
  \label{thin:eq-active-action-layers}
\end{equation}
At a wall \(x_n=T\), this gives \(J_{\mathrm{rad}}=n-1\), so the equality
layer is excluded.

For every included \(x_n<T\), let \(x=R(x_n)>0\).  The exact count
\eqref{thin:eq-half-integer-count} satisfies
\[
  M_\varepsilon(x)<\frac{x}{\varepsilon}+\frac12.
\]
The two sides are nonnegative, so after squaring and using \(x\leq a\),
\begin{equation}
  \varepsilon^2M_\varepsilon(x)^2
  <x^2+\varepsilon x+\frac{\varepsilon^2}{4}
  \leq x^2+\varepsilon a+\frac{\varepsilon^2}{4}.
  \label{thin:eq-angular-rounding}
\end{equation}
At a half-integer wall \(x/\varepsilon=m+1/2\), the strict count is \(m\),
so \eqref{thin:eq-angular-rounding} remains strict.  Summing and using
\eqref{thin:eq-active-action-layers} gives
\begin{equation}
  \varepsilon^2
  \sum_{\substack{n\geq1\\x_n<T}}
    M_\varepsilon(R(x_n))^2
  <\sum_{\substack{n\geq1\\x_n<T}}F(x_n)+E,
  \label{thin:eq-complete-angular-error}
\end{equation}
where
\begin{equation}
  E=\left(\frac a\pi+\frac14\right)
     \left(\varepsilon a+\frac{\varepsilon^2}{4}\right).
  \label{thin:eq-angular-error}
\end{equation}

By Lemma~\ref{thin:lem-radial-deficit} and
\(a\geq\pi/(4\varepsilon)\),
\begin{equation}
  \frac{D_{\mathrm{rad}}}{a^2}
  \geq\frac{\rho^2}{4}-\frac{\pi\rho}{4a}
  \geq\frac{(1-\varepsilon)(1-5\varepsilon)}4.
  \label{thin:eq-high-lower-screen}
\end{equation}
The last expression has derivative
\((-6+10\varepsilon)/4<0\) on \(0<\varepsilon\leq1/8\), and therefore
\begin{equation}
  \frac{D_{\mathrm{rad}}}{a^2}\geq\frac{21}{256}.
  \label{thin:eq-high-lower-endpoint}
\end{equation}
On the other hand, expansion of \eqref{thin:eq-angular-error} and
\(1/a\leq4\varepsilon/\pi\) give
\begin{equation}
  \frac{E}{a^2}
  \leq\frac\varepsilon\pi+\frac{\varepsilon^2}\pi
     +\frac{\varepsilon^3}{\pi^2}
     +\frac{\varepsilon^4}{\pi^2}.
  \label{thin:eq-high-upper-screen}
\end{equation}
This scaling is deliberately non-strict: equality occurs at the common
optical face \(a=\pi/(4\varepsilon)\).  The right-hand side is strictly
increasing in \(\varepsilon>0\).  At \(\varepsilon=1/8\), the elementary
strict inequality \(\pi>3\) gives
\begin{equation}
  \frac{E}{a^2}
  \leq\frac1{8\pi}+\frac1{64\pi}
      +\frac1{512\pi^2}+\frac1{4096\pi^2}
  <\frac{193}{4096}.
  \label{thin:eq-high-upper-endpoint}
\end{equation}
Combining \eqref{thin:eq-high-lower-endpoint} and
\eqref{thin:eq-high-upper-endpoint} gives the exact positive reserve
\begin{equation}
  D_{\mathrm{rad}}-E
  >a^2\left(\frac{21}{256}-\frac{193}{4096}\right)
  =\frac{143}{4096}a^2>0.
  \label{thin:eq-high-reserve}
\end{equation}

Equations \eqref{thin:eq-radial-deficit-definition},
\eqref{thin:eq-complete-angular-error}, and
\eqref{thin:eq-high-reserve} now give the complete strict chain
\begin{equation}
  \varepsilon^2
  \sum_{\substack{n\geq1\\x_n<T}}
    M_\varepsilon(R(x_n))^2
  <\sum_{\substack{n\geq1\\x_n<T}}F(x_n)+E
  <\int_0^T F(t)\,dt.
  \label{thin:eq-high-strict-chain}
\end{equation}
By Lemma~\ref{thin:lem-layer-cake}, this is \(P_{\mathcal A}<I\), and hence
\[
  N_D(A_{\rho,1},K^2)
  \leq\frac{2P_{\mathcal A}}{\varepsilon^2}
  <\frac{2I}{\varepsilon^2}
  =\frac{2}{9\pi}(1-\rho^3)K^3.
\]

Every inequality used above includes \(a=\pi/(4\varepsilon)\).  At that
face the non-strict scalings in
\eqref{thin:eq-high-lower-screen}--\eqref{thin:eq-high-upper-screen} are
equalities, but the strict estimate \(\pi>3\) retains the reserve
\eqref{thin:eq-high-reserve}.  Thus the high proof independently owns the
whole common face, including \((\varepsilon,a)=(1/8,2\pi)\).
\end{proof}

\begin{lemma}[Thin-shell endpoint]
\label{lem:thin-shell}
For \(7/8\leq\rho<1\) and \(K\geq0\),
\begin{equation}
  N_D(A_{\rho,1},K^2)
  \leq\frac{2}{9\pi}(1-\rho^3)K^3.
  \label{eq:thin-shell-theorem}
\end{equation}
Equality holds at \(K=0\), and the displayed inequality is strict when
\(K>0\).
\end{lemma}

\begin{proof}
At \(K=0\), positivity of the Dirichlet operator gives a zero strict count,
and the right-hand side is also zero.  Suppose \(K>0\).  Then
\(0<\varepsilon\leq1/8\) and \(a>0\).  The two inclusive intervals
\[
  0\leq a\leq\frac{\pi}{4\varepsilon},
  \qquad
  a\geq\frac{\pi}{4\varepsilon}
\]
cover every \(a\geq0\), with no gap and no third range.  The first is proved
by Proposition~\ref{thin:prop-low-piece}; the second is proved by
Proposition~\ref{thin:prop-high-piece}.  Both propositions include their
common face.  Formula \eqref{thin:eq-weyl-normalization} identifies the
low-piece target with the right-hand side of
\eqref{eq:thin-shell-theorem}.  This completes the proof on the closed face
\(\rho=7/8\) and for every \(7/8<\rho<1\).  The limit
\(\rho\uparrow1\) means \(\varepsilon\downarrow0\) through positive values;
the nonexistent shell \(\varepsilon=0\) is never inserted into a
denominator or treated as part of the domain.
\end{proof}

\begin{remark}[Dependency boundary]
\label{thin:rem-no-lorch}
No finite list of Bessel zeros, no numerical root certificate, and no lower
bound for a first Bessel zero is used in the proof of
Lemma~\ref{lem:thin-shell}.  In particular, the inequalities of Lorch for
\(j_{\nu,1}\) are dependencies of the compact-middle finite-channel
staircase, not of the endpoint \(7/8\leq\rho<1\).  The only non-elementary
input specific to the high piece is the all-domain optimized phase estimate
\eqref{thin:eq-phase-input}, whose proof appears earlier in this manuscript;
all product, action, inverse, curvature, Stieltjes, rounding, seam, and
constant estimates used after that input have been proved above.
\end{remark}

% Self-contained compact-middle proof fragment.
% It assumes only the notation and phase-to-lattice reduction established
% earlier in the manuscript.

\section{The compact-middle argument in full detail}
\label{cm:sec-full}

We retain the strict counting convention
\([x]_<:=\#\{n\in\mathbb N:n<x\}\), put
\[
 W(\rho,K):=\frac{2}{9\pi}(1-\rho^3)K^3,
 \qquad z_\rho:=\frac{\pi}{1-\rho},
 \qquad k_m(\rho):=\sqrt{z_\rho^2+m(m+1)},
\]
and denote by \(q_{\ell,n}(\rho)\) the \(n\)-th positive frequency in
angular channel \(\ell\).  Its multiplicity is \(2\ell+1\).  The phase
reduction proved earlier gives
\begin{equation}
 N_D(A_{\rho,1},K^2)
 \le P_{\rm coarse}(\rho,K)
 :=\sum_{\ell+1/2<K}(2\ell+1)
 \left[A_{\rho,K}(\ell+1/2)+\frac14\right]_<,
 \label{cm:eq-coarse}
\end{equation}
where \(A_{\rho,K}=G_K-G_{\rho K}\) and
\[
 G_a(\nu)=
 \begin{cases}
 \displaystyle\frac1\pi\left(\sqrt{a^2-\nu^2}
 -\nu\arccos\frac\nu a\right),&0\leq\nu<a,\\[4pt]
 0,&\nu\geq a.
 \end{cases}
\]
The value at both interfaces is therefore fixed by zero extension.  This
will matter at \(\nu=\rho K\), at \(\nu=K\), and at every integer proxy
wall.

\subsection{Constants and the primitive analytic envelope}

For \(0\leq s\leq1\), set
\[
 G_1(s)=\frac1\pi\left(\sqrt{1-s^2}-s\arccos s\right),
 \quad
 \omega_0=G_1(1/2)=\frac{\sqrt3}{2\pi}-\frac16,
\]
\[
 C_*=\frac12+\frac{8\sqrt2}{15},\qquad
 C_0=1+\frac{8\sqrt2}{15},\qquad
 \rho_*:=\frac{\omega_0}{1+2C_*}=\frac{\omega_0}{2C_0}.
\]
For \(0<\rho<1\), define
\[
 a_\rho=\frac{2\pi\rho}{1-\rho},\qquad
 \eta_\rho=G_1\!\left(\max\{\rho,1/2\}\right),
\]
\begin{equation}
 K_0(\rho)=
 \left(
 \frac{\sqrt{a_\rho}+\sqrt{a_\rho+4\eta_\rho C_0}}
 {2\eta_\rho}
 \right)^2,
 \qquad
 H_0(\rho)=\frac{C_*}{\omega_0-\rho}\quad(\rho<\omega_0).
 \label{cm:eq-upper-functions}
\end{equation}
We first re-establish every primitive owner that is needed to form the
compact residual.  This prevents the residual description from hiding an
unproved coverage assumption.

\begin{proposition}[Primitive analytic owners]
\label{cm:prop-primitive-owners}
The P\'olya bound holds on each of the following inclusive domains:
\begin{enumerate}
\item \(2\rho K\leq1\);
\item \((1-\rho)K\leq\pi\);
\item \(0<\rho<\omega_0\) and
      \(K(\omega_0-\rho)\geq C_*\), equivalently \(K\geq H_0(\rho)\);
\item \(0<\rho\leq\rho_*\), for every \(K\geq0\);
\item \(K\geq K_0(\rho)\);
\item \(0<1-\rho\leq1/6\) and
      \(K(1-\rho)^2\geq54\);
\item \(7/8\leq\rho<1\), for every \(K\geq0\).
\end{enumerate}
In particular, on the full ratio interval \(0<\rho<1\), the bound holds
whenever \(K\geq295^2\).
\end{proposition}

\begin{proof}
For item 1, \(\mu=\rho K\leq1/2\); every shifted half-integer tail begins
at or above the inner interface and is controlled by the convex high-angular
tail inequality.  Summing the tail integrals with the weighted identity
gives \(N_D\leq W\).  For item 2,
\(K\leq z_\rho=q_{0,1}\), and
\eqref{cm:eq-product-bound} (whose elementary proof is repeated below)
shows that no radial mode is strictly below \(K\).

We record the tail bookkeeping used in items 3--5.  Put
\[
 x_r=r+\frac12,\qquad
 q_\ell=\left\lfloor A_{\rho,K}\!\left(\ell+\frac12\right)
                  +\frac14\right\rfloor
\]
and
\[
 \mathcal T_r=
 \left\lfloor A_{\rho,K}(x_r)+\frac14\right\rfloor
 +2\sum_{j\geq1}
 \left\lfloor A_{\rho,K}(x_r+j)+\frac14\right\rfloor .
\]
All sums are finite because \(A_{\rho,K}=0\) at and beyond \(K\).
For integers \(a<b\), write
\[
 \mathbf T(h,a,b)=\frac12\lfloor h(a)\rfloor
 +\sum_{j=a+1}^{b-1}\lfloor h(j)\rfloor
 +\frac12\lfloor h(b)\rfloor .
\]
The elementary multiplicity identity
\[
 \sum_{\ell\geq0}(2\ell+1)q_\ell
 =\sum_{r\geq0}\left(q_r+2\sum_{\ell>r}q_\ell\right)
 =\sum_{r\geq0}\mathcal T_r
\]
shows that \(P_{\rm coarse}\leq\sum_r\mathcal T_r\).  If every active
tail satisfies
\[
 \mathcal T_r\leq2\int_{x_r}^{K}A_{\rho,K}(z)\,dz,
\]
then, with
\(f_3(z)=\#\{r\geq0:x_r\leq z\}\) and
\(F_3(z)=\int_0^zf_3(t)\,dt\),
\[
 P_{\rm coarse}\leq2\int_0^K f_3(z)A_{\rho,K}(z)\,dz.
\]
Writing \(z=m+1/2+s\), \(m\in\mathbb N_0\), \(0\leq s<1\), gives
\[
 \frac{z^2}{2}-F_3(z)=\frac12\left(s-\frac12\right)^2\geq0;
\]
for \(0\leq z<1/2\), the same inequality is immediate.  Since
\(A_{\rho,K}\) is decreasing, absolutely continuous, and vanishes at
\(K\), integration by parts yields
\[
 P_{\rm coarse}
 \leq-\int_0^Kz^2A_{\rho,K}'(z)\,dz
 =\int_0^K2zA_{\rho,K}(z)\,dz
 =\frac{2}{9\pi}(1-\rho^3)K^3=W(\rho,K).
\]
For the final equality, scaling reduces each \(G_a\)-integral to
\[
 \int_0^1 2xG_1(x)\,dx
 =\frac2\pi\left(\frac13-\frac29\right)=\frac{2}{9\pi}.
\]

We shall use the exact trapezoidal-floor estimates proved in the
one-dimensional reduction.  For a nonnegative concave \(h\),
\(\mathbf T(h,a,b)\leq\int_a^bh(t)\,dt\).  If \(g\) is nonnegative,
decreasing, convex, and \(1/2\)-Lipschitz, vanishes at \(b\), and is
\(1/3\)-Lipschitz after \(t_*\), then
\[
 \mathbf T(g+1/4,a,b)
 \leq\int_a^bg(t)\,dt-\frac14\lfloor g(t_*)\rfloor.
\]
The corresponding shifted convex estimate gives
\(\mathcal T_r\leq2\int_{x_r}^KG_K(z)\,dz\) whenever \(x_r\geq\mu\).
Indeed, on that interval \(A=G_K\), and
\[
 G_K'(z)=-\frac1\pi\arccos\frac zK\in[-1/2,0],
 \qquad
 G_K''(z)=\frac1{\pi\sqrt{K^2-z^2}}>0.
\]

We now prove item 3.  Write \(A=A_{\rho,K}\).  Only a start
\(x_0<\mu=\rho K\) remains.  Put
\[
 n=\lfloor\mu-x_0\rfloor,\qquad q=x_0+n,\qquad
 B=\lceil K-x_0\rceil .
\]
Thus \(q\leq\mu<q+1\), \(n<B\), and zero extension gives
\(G_K(x_0+B)=0\).  The action is concave on \([x_0,q]\), because
\[
 A''(z)=\frac1{\pi\sqrt{K^2-z^2}}
 -\frac1{\pi\sqrt{\mu^2-z^2}}<0\qquad(0<z<\mu).
\]
Moreover \(q\leq\mu=\rho K<\omega_0K<K/2\).  Splitting at \(n\), using
the concave estimate on the head and the improved convex estimate on the
\(G_K\)-suffix with \(t_*=K/2-x_0\), gives, also when \(n=0\),
\[
 \frac{\mathcal T_r}{2}
 \leq\int_{x_0}^{q}A(z)\,dz+\frac n4
 +\int_q^KG_K(z)\,dz-\frac14\lfloor\omega_0K\rfloor.
\]
The only integral mismatch is
\[
 \delta:=\int_q^\mu G_\mu(z)\,dz.
\]
To bound it, write \(t=\cos\theta\).  Concavity of \(\sin\) on
\([0,\pi/4]\) gives
\(\sin(\theta/2)\geq\sqrt2\,\theta/\pi\), hence
\(\arccos t\leq(\pi/2)\sqrt{1-t}\) for \(0\leq t\leq1\).  Since
\[
 G_\mu(z)=\frac1\pi\int_z^\mu\arccos\frac t\mu\,dt,
\]
we obtain
\[
 0\leq\delta
 \leq\frac{2(\mu-q)^{5/2}}{15\sqrt\mu}
 <\frac{2\sqrt2}{15},
\]
where a low tail has \(\mu>x_0\geq1/2\) and \(0\leq\mu-q<1\).
Consequently
\[
 \frac{\mathcal T_r}{2}
 \leq\int_{x_0}^{K}A(z)\,dz
 +\delta-\frac{\lfloor\omega_0K\rfloor-n}{4}.
\]
Finally \(n\leq\rho K-1/2\) and \(\lfloor y\rfloor>y-1\), so
\[
 \lfloor\omega_0K\rfloor-n
 >K(\omega_0-\rho)-\frac12
 \geq C_*-\frac12=\frac{8\sqrt2}{15}>4\delta.
\]
Thus every low-interface tail is strict, every high-interface tail is
non-strict, and the weighted scaffold proves item 3, including the face
\(K(\omega_0-\rho)=C_*\).

For item 4, if \(\mu\leq1/2\) use item 1.  Otherwise \(K>1/(2\rho)\).
Because \(\rho_*=\omega_0/(1+2C_*)\), for
\(0<\rho\leq\rho_*\),
\[
 K(\omega_0-\rho)>
 \frac{\omega_0-\rho}{2\rho}
 \geq\frac{\omega_0-\rho_*}{2\rho_*}=C_*,
\]
and item 3 applies.  All faces are included; at \(K=0\) both sides vanish.

For item 5, fix a low-interface tail \(x_0=r+1/2<\mu\), retain the
preceding definition of \(\mathcal T_r\), and write \(A=A_{\rho,K}\).
Put \(n=\lfloor\mu-x_0\rfloor\), \(q=x_0+n\), and
\(B=\lceil K-x_0\rceil\).  On \([0,n]\),
\(h(t)=A(x_0+t)+1/4\) is decreasing, concave, and Lipschitz with
constant \(c_\rho=\arccos\rho/\pi\).  If \(p\) is the last index of its
first constant ordinary-floor shelf (put \(p=n\) when the shelf never
drops), the concave trapezoidal-floor inequality gives
\begin{equation}
 \mathbf T(h,0,n)
 \leq\int_0^nA(x_0+t)\,dt+\frac n4
 -\frac{1-c_\rho}{2}(n-p).
 \label{cm:eq-fixed-head}
\end{equation}
For \(n=0\), the head is absent and \(p=0\).

Let \(\vartheta=-A'\).  Direct differentiation gives
\[
 \vartheta'(z)=\frac1{\pi\sqrt{\mu^2-z^2}}
 -\frac1{\pi\sqrt{K^2-z^2}}
 \geq\frac{1-\rho}{\pi\rho K}=:\kappa_\rho.
\]
Equal floors on the first shelf imply
\(0\leq A(x_0)-A(x_0+p)<1\), including at an integer endpoint.  Hence
\[
 1>\int_{x_0}^{x_0+p}\vartheta(z)\,dz
 \geq\frac{\kappa_\rho p^2}{2},
 \qquad p<\sqrt{a_\rho K}.
\]

For the convex tail set \(s=\max\{q,K/2\}\).  The function
\(G_K(x_0+t)\) is decreasing and convex, is Lipschitz \(1/2\) on the
whole tail and \(1/3\) after \(s\), and
\(G_K(s)\geq K\eta_\rho\).  The convex trapezoidal-floor inequality
therefore gives
\begin{equation}
 \mathbf T(G_K(x_0+\cdot)+1/4,n,B)
 \leq\int_n^B G_K(x_0+t)\,dt
 -\frac14\lfloor K\eta_\rho\rfloor.
 \label{cm:eq-fixed-tail}
\end{equation}
At \(j=n\), the two half endpoint weights in
\eqref{cm:eq-fixed-head}--\eqref{cm:eq-fixed-tail} majorize the original
full \(A\)-sample.  Their only integral mismatch is
\[
 \delta=\int_q^\mu G_\mu(z)\,dz,
 \qquad 0\leq\delta\leq
 \frac{2(\mu-q)^{5/2}}{15\sqrt\mu}<\frac{2\sqrt2}{15}.
\]
Writing \(d_\rho=1-2c_\rho=2\arcsin\rho/\pi\), combination of the
head and tail yields
\[
 \frac{\mathcal T_r}{2}
 \leq\int_{x_0}^KA(z)\,dz+\delta-\frac M4,
 \quad
 M=\lfloor K\eta_\rho\rfloor+d_\rho(n-p)-p.
\]
Since \(d_\rho>0\) and \(p<\sqrt{a_\rho K}\), the strict floor
inequality, valid even when \(K\eta_\rho\) is an integer, gives
\[
 M>K\eta_\rho-1-\sqrt{a_\rho K}.
\]
Thus the exact sufficient condition is
\begin{equation}
 K\eta_\rho-\sqrt{a_\rho K}\geq C_0.
 \label{cm:eq-fixed-ratio-condition}
\end{equation}
Indeed it implies \(M>8\sqrt2/15>4\delta\), so every low-interface tail
is strictly below twice its action integral.  Tails beginning at or above
\(\mu\) satisfy the non-strict convex high-interface bound; the weighted
identity completes the proof.  The cases \(n=0\), \(p=0\), \(p=n\),
\(q=\mu\), \(q\) on either side of \(K/2\), and all strict-floor walls
are included in the displayed inequalities.

Writing \(Y=\sqrt K\), the left side minus \(C_0\) is
\(\eta_\rho Y^2-\sqrt{a_\rho}Y-C_0\).  It has one positive root, and
the square of that root is precisely \(K_0(\rho)\) in
\eqref{cm:eq-upper-functions}.  Hence \eqref{cm:eq-fixed-ratio-condition}
holds, including equality, exactly for \(K\geq K_0(\rho)\).

We give the seam proof for item 6 because it is the nontrivial owner in the
primitive mask.  Put
\[
 \varepsilon=1-\rho,\quad y=\sqrt\varepsilon,\quad
 \kappa=K\varepsilon^2,\quad \mu=\rho K,\quad
 c=\frac{\arccos\rho}{\pi},\quad d=1-2c,\quad
 \eta=G_1(\rho).
\]
For a shifted low-interface tail beginning at
\(x_0=r_0+1/2<\mu\), write
\[
 n=\lfloor\mu-x_0\rfloor,\quad q=x_0+n,\quad
 \beta=\mu-q,\quad
 h_j=\left\lfloor A(x_0+j)+\frac14\right\rfloor.
\]
Let \(p\) be the last index of the first constant plateau, set
\(m=n-p\), \(R=p-dm\), and use \(p=m=R=0\) if \(n=0\).  The exact
one-tail decomposition is
\begin{equation}
 \frac{\mathcal T_{r_0}}2
 \leq\int_{x_0}^K A(z)\,dz+\delta-\frac M4,
 \quad M=\lfloor K\eta\rfloor-R,
 \quad 0\leq\delta<\frac{2\sqrt2}{15}.
 \label{cm:eq-seam-tail}
\end{equation}
The curvature shelf estimate gives
\begin{equation}
 p<\sqrt{\frac{2\pi\rho K}{\varepsilon}}.
 \label{cm:eq-seam-shelf}
\end{equation}
It remains to prove \(M>4\delta\) on
\(0<y\leq1/\sqrt6\), \(\kappa\geq54\).

Set \(P=py\), \(r=Ry\), and \(S=(p+m)y\).  The elementary half-angle
estimate \(\cos(3\pi/16)<5/6\leq\rho\) gives
\begin{equation}
 d>5/8,\qquad R>0\Longrightarrow
 m<\frac85p,\quad S<\frac{13}{5}P.
 \label{cm:eq-seam-angular}
\end{equation}
With \(t=x_0/(\rho K)\), \(w=1-t\), equations
\eqref{cm:eq-seam-shelf}--\eqref{cm:eq-seam-angular} imply
\[
 w<\widehat q
 :=\frac{y^4+y^3S}{\kappa\rho}
 <\frac1{1620}+\frac{13}{5}\frac{763}{5000}
 =\frac{804689}{2025000}<\frac{159}{400}.
\]
The estimates use only
\(\pi<1571/500\),
\(140/99<\sqrt2<99/70\), and
\(y<49/120\), each proved by positive rational squaring.  In particular,
\(x_0/K>241/480>1/2\).

For \(0<z<\mu\), let \(\sigma(z)=-A'(z)\).  Direct differentiation gives
\[
 \sigma'(z)>0,\qquad
 \sigma(x_0)\geq
 \frac{\varepsilon t}{\pi\sqrt{1-\rho^2t^2}}.
\]
Equal ordinary floors along the plateau imply the strict inequality
\(A(x_0)-A(x_0+p)<1\), even at an integer wall.  Put
\[
 Q=1+\frac{y^2}{\kappa}+\frac{yS}{\kappa}.
\]
Since \(1-\rho^2t^2<2y^2Q\), integration of \(\sigma\) across the
plateau yields
\begin{equation}
 P^2<\frac{2\pi^2Q}{(1-\widehat q)^2}.
 \label{cm:eq-seam-selfconsistent}
\end{equation}
Using \(d>5/8\), replace
\(S\) by \(S_*=(13P-8r)/5\), and define
\[
 Q_*=1+\frac{y^2}{\kappa}+\frac{yS_*}{\kappa},
 \quad q_*=\frac{y^2}{\rho}(Q_*-1),
 \quad H(P,r)=\frac{2\pi^2Q_*}{(1-q_*)^2}.
\]
The right side increases with \(Q_*\), so
\(P^2<H(P,r)\).  Suppose \(r\geq B:=14/3\).  Since
\(R=p-dm\leq p\), one has \(0<r\leq P\).  At fixed \(P\), the quantity
\(S_*(P,r)=(13P-8r)/5\) decreases with \(r\); hence \(Q_*\) and \(q_*\)
also decrease.  Because \(q_*<159/400<1\), the expression
\(2\pi^2Q_*/(1-q_*)^2\) increases with \(Q_*\).  Therefore
\[
 H(P,r)\leq H(P,B).
\]
Now freeze \(r=B\) and follow the entire segment \(B\leq P'\leq P\).
The original shelf ceiling remains valid because \(P'\leq P\), while
\[
 0<S_*(P',B)\leq S_*(P',0)=\frac{13}{5}P'.
\]
Thus the bounds \(q_*<159/400\) and all positive denominators used above
remain valid along the path.  On it,
\[
 \partial_{P'}H(P',B)<\frac{16}{5},\qquad
 \frac{d}{dP'}\bigl(P'^2-H(P',B)\bigr)>\frac{92}{15},
\]
whereas the initial exact reserve is
\[
 B^2-H(B,B)>
 \frac{2505132463469}{2616573970125}>0.
\]
Integration of the positive derivative yields
\[
 P^2-H(P,B)>B^2-H(B,B)>0.
\]
Consequently \(P^2>H(P,B)\geq H(P,r)\), contradicting
\(P^2<H(P,r)\) from \eqref{cm:eq-seam-selfconsistent}.  Hence
\begin{equation}
 R<\frac{14}{3\sqrt\varepsilon}.
 \label{cm:eq-seam-R}
\end{equation}

Finally
\[
 \eta=\frac1\pi\int_0^\varepsilon\arccos(1-v)\,dv
 \geq\frac{2\sqrt2}{3\pi}\varepsilon^{3/2},
\]
so on \(\kappa\geq54\),
\[
 K\eta>\frac{280000}{17281}\frac1y.
\]
If \(R>0\), the strict inequality \(\lfloor x\rfloor>x-1\),
\eqref{cm:eq-seam-R}, and \(1/y>120/49\) give
\[
 M-4\delta>
 \left(\frac{598066}{51843}\right)\frac{120}{49}
 -1-\frac{132}{175}
 =\frac{80132733}{3024175}>0.
\]
If \(R\leq0\), the stronger reserve is
\(114694733/3024175>0\).  Substitution in
\eqref{cm:eq-seam-tail} makes every low-interface tail strictly smaller
than its action integral; high-interface tails are controlled by convexity.
The weighted identity then proves item 6.  This argument includes
\(\kappa=54\), every plateau wall, \(R=0\), every integer
\(K\eta\) wall, and the interface \(x_0=\mu\).
The derivative cap, initial reserve, and terminal rational payments in this
synthetic comparison are regenerated line by line by the ledger printed in
  Parts~I--VI of the accompanying \emph{Purely Analytic Supplement}.

Item 7 is the two-piece thin-shell theorem proved by the exact product
deficit for low optical frequency and the complementary-action Stieltjes
deficit for high optical frequency.  For completeness, its two closing
reserves are
\[
 \frac25-\left(\frac16+\frac1{32}+\frac1{576}\right)
 =\frac{577}{2880}>0,
 \qquad
 D_{\rm rad}-E>\frac{143}{4096}a^2>0,
\]
with \(a=(1-\rho)K\); the common face is included.  Thus it supplies the
entire closed ratio face \(\rho=7/8\) needed here.

It remains to prove the global frequency stated in the proposition.  On
\([\rho_*,1/16]\), item 3 and the increase of \(H_0\) leave only
\(K<64\).  On
\([1/16,5/6]\), monotonicity of \(K_0\), proved below, and the exact
estimate \(K_0(5/6)<295^2\) leave only \(K<295^2\).  On
\([5/6,7/8]\), item 6 leaves only
\(K<54/(1-\rho)^2\leq3456\).  Item 7 owns every frequency on
\([7/8,1)\).  All threshold faces are included, so \(K=295^2\) is
covered at every ratio.
\end{proof}

Proposition~\ref{cm:prop-primitive-owners} covers the following inclusive
sets in \([\rho_*,7/8]\times[0,\infty)\):
\begin{align*}
 &\rho=\rho_*,\quad \rho=7/8,\quad 2\rho K\leq1,
 \quad (1-\rho)K\leq\pi,\\
 &\rho<\omega_0,\quad K(\omega_0-\rho)\geq C_*,
 \qquad K\geq K_0(\rho),\\
 &0<1-\rho\leq1/6,\quad K(1-\rho)^2\geq54,
 \qquad K\geq295^2.
\end{align*}

Put
\[
 L(\rho)=\max\left\{\frac1{2\rho},\frac\pi{1-\rho}\right\},
\]
and let \(U(\rho)\) be the minimum of all eligible upper owners in
\eqref{cm:eq-upper-functions} and the seam owner \(54/(1-\rho)^2\).
Since every primitive frequency wall is included, exact complementation
gives
\begin{equation}
 \mathcal D_{16}=
 \left\{(\rho,K):\rho_*<\rho<\frac78,\quad
 L(\rho)<K<U(\rho)\right\}.
 \label{cm:eq-D16}
\end{equation}
Both frequency inequalities are strict.

We record the switches, since they determine every later face owner.  The
lower walls meet at
\[
 \rho_c:=\frac1{1+2\pi},\qquad
 L(\rho)=\begin{cases}(2\rho)^{-1},&\rho\leq\rho_c,\\
 z_\rho,&\rho\geq\rho_c.
 \end{cases}
\]
On \((\rho_*,\omega_0)\), substituting \(K=H_0(\rho)\) into the
positive-root equation
\(\eta_\rho K-\sqrt{a_\rho K}-C_0=0\) gives, after a reversible
squaring,
\begin{equation}
 P(\rho):=(1-\rho)\left(C_0\rho-\frac{\omega_0}{2}\right)^2
 -2\pi C_*\rho(\omega_0-\rho)=0.
 \label{cm:eq-Pswitch}
\end{equation}
Here \(C_0\rho-\omega_0/2>0\), and
\[
 P''(\rho)=2C_0^2(1-3\rho)+2C_0\omega_0+4\pi C_*>0.
\]
At \(\rho_*\) the squared term vanishes and the other term is negative,
whereas
\(P(\omega_0)=(1-\omega_0)\omega_0^2C_*^2>0\).  Strict convexity
therefore gives a unique switch \(\rho_{HK}\).  Here is a fully rational
witness for its locator.  Put
\[
 S_n(x)=\sum_{j=0}^{n-1}\frac{(-1)^jx^{2j+1}}{2j+1}.
\]
The alternating series in Machin's identity gives
\[
\begin{aligned}
 \frac{15707963}{5000000}
 &<16S_6(1/5)-4S_3(1/239)<\pi\\
 &<16S_5(1/5)-4S_2(1/239)
 <\frac{31415927}{10000000}.
\end{aligned}
\]
The two outer gaps are respectively
\[
 \frac{88855616195735287}{1688778771196584375000000},
 \qquad
 \frac{3783145391}{215017724250000000}.
\]
Set these outer bounds equal to \(\pi^-\) and \(\pi^+\), and set
\[
\begin{array}{ll}
 s_2^-=\dfrac{2828427}{2000000},&
 s_2^+=\dfrac{1767767}{1250000},\\[3pt]
 s_3^-=\dfrac{4330127}{2500000},&
 s_3^+=\dfrac{17320509}{10000000}.
\end{array}
\]
Their square-test residuals are
\[
\begin{array}{ll}
2(2000000)^2-2828427^2=705671,&
1767767^2-2(1250000)^2=166289,\\
3(2500000)^2-4330127^2=163871,&
17320509^2-3(10000000)^2=32019081,
\end{array}
\]
so \(s_j^-<\sqrt j<s_j^+\).  Define
\[
 \omega^-=\frac{s_3^-}{2\pi^+}-\frac16,\quad
 \omega^+=\frac{s_3^+}{2\pi^-}-\frac16,\quad
 C_*^\pm=\frac12+\frac8{15}s_2^\pm,\quad
 C_0^\pm=1+\frac8{15}s_2^\pm .
\]
For \(r_-=9407/100000\), monotonic substitution gives
\[
 P(r_-)\leq\mathcal P_-:=
 (1-r_-)(C_0^+r_--\omega^-/2)^2
 -2\pi^-C_*^-r_-(\omega^--r_-),
\]
and exact reduction gives
\[
 \mathcal P_-+\frac1{5000000}
 =-\frac{14699034939621959968377942867507773}
 {2409571458177072753906250000000000000000000}<0.
\]
For \(r_+=94071/1000000\), one first checks
\[
 C_0^-r_+-\frac{\omega^+}{2}-\frac1{10}
 =\frac{619967074526802271}
 {58904861250000000000}>0.
\]
Thus the square has the same monotonic direction, and
\[
 P(r_+)\geq\mathcal P_+:=
 (1-r_+)(C_0^-r_+-\omega^+/2)^2
 -2\pi^+C_*^+r_+(\omega^+-r_+),
\]
where
\[
 \mathcal P_+-\frac7{10000000}
 =\frac{118239888922618787414317706992300967689}
 {3469782678881751562500000000000000000000000000}>0.
\]
Consequently
\[
 \frac{9407}{100000}<\rho_{HK}<\frac{94071}{1000000},
\]
such that \(H_0\) is active below it and \(K_0\) above it.  The locator
is not used as a replacement for the implicit root.

The function \(K_0\) is increasing: \(a_\rho\) increases,
\(\eta_\rho\) is constant up to \(1/2\) and then decreases because
\(G_1'(s)=-\arccos(s)/\pi\), and the positive root of
\(\eta y^2-\sqrt a\,y-C_0=0\) increases with \(a\) and decreases with
\(\eta\).  At \(\rho=5/6\),
\[
 \cos\frac9{10}
 <1-\frac{(9/10)^2}{2}+\frac{(9/10)^4}{24}
 =\frac{49787}{80000}<\frac56,
\]
so \(\arccos(5/6)<9/10\).  Since
\[
 \eta_{5/6}=\frac1\pi\int_{5/6}^1\arccos t\,dt
 <\frac1{6\pi}\arccos\frac56<\frac1{20},
\]
where \(\pi>3\), the positive-root equation implies
\[
 K_0(5/6)>\frac{a_{5/6}}{\eta_{5/6}^2}
 >10\cdot3\cdot20^2=12000,
 \qquad
 \frac{54}{(1-\rho)^2}\leq3456
 \quad(5/6\leq\rho<7/8).
\]
Consequently
\begin{equation}
 U(\rho)=
 \begin{cases}
 H_0(\rho),&\rho_*<\rho<\rho_{HK},\\
 K_0(\rho),&\rho_{HK}\leq\rho<5/6,\\
 54/(1-\rho)^2,&5/6\leq\rho<7/8.
 \end{cases}
 \label{cm:eq-U-piecewise}
\end{equation}
The values agree at \(\rho_{HK}\); the seam face \(\rho=5/6\) is
owned by the third branch.  For the opposite estimate, the inequality
\(\arccos t\geq\sqrt{2(1-t)}\), obtained from
\(\cos u\geq1-u^2/2\), gives
\[
 \eta_{5/6}\geq\frac1\pi\int_{5/6}^1\sqrt{2(1-t)}\,dt
 =\frac1{9\pi\sqrt3}>\frac1{49}.
\]
Indeed, \(\pi<22/7\), while
\(343^2-3\cdot198^2=37>0\), so
\(9\pi\sqrt3<9(22/7)(343/198)=49\).
Also
\[
 a_{5/6}=10\pi<\frac{220}{7}<36,\qquad
 C_0<1+\frac8{15}\frac{99}{70}=\frac{307}{175},
\]
where \(99^2-2\cdot70^2=1>0\).  Therefore, at \(Y=295\),
\[
 \eta_{5/6}Y^2-\sqrt{a_{5/6}}Y-C_0
 >\frac{295^2}{49}-6\cdot295-\frac{307}{175}
 =\frac{5226}{1225}>0,
\]
so \(K_0(5/6)<295^2\).

For the remaining \(H_0\) check, \(\sqrt3>12/7\) and \(\pi<22/7\)
give \(15\sqrt3>8\pi\), hence \(\omega_0>1/10\).  Also
\(\sqrt2<3/2\) gives \(C_*<13/10\).  Consequently
\[
 \omega_0-\frac1{16}>\frac3{80}>\frac{19}{528},
 \qquad
 H_0(1/16)<\frac{13/10}{19/528}
 =\frac{3432}{95}<64.
\]
Together with the seam ceiling \(3456\), this proves that
\(K=295^2\) is covered at every ratio, including equality.

\subsection{Variational comparisons and the complete zero registry}

The transformed channel form is
\[
 \mathfrak a_\ell[u]=\int_\rho^1
 \left(|u'|^2+\frac{\ell(\ell+1)}{r^2}|u|^2\right)dr,
 \qquad D(\mathfrak a_\ell)=H_0^1(\rho,1).
\]
Since \(r^{-2}\geq1\), the min--max principle gives, at every radial
index,
\begin{equation}
 q_{\ell,n}^2\geq n^2z_\rho^2+\ell(\ell+1),
 \qquad q_{0,n}=nz_\rho.
 \label{cm:eq-product-bound}
\end{equation}
Extension by zero from \((\rho,1)\) to \((0,1)\) embeds the shell
fixed-channel form domain isometrically into the ball fixed-channel form
domain.  The zero trace at \(r=\rho\) preserves membership in
\(H_0^1(0,1)\), and Hardy's inequality controls the centrifugal term.
Thus, preserving both \(\ell\) and \(n\),
\begin{equation}
 q_{\ell,n}^2\geq j_{\ell+1/2,n}^2.
 \label{cm:eq-shell-ball}
\end{equation}
The ball forms have the same domain for every angular order; hence, for
\(p\geq\ell\),
\begin{equation}
 j_{p+1/2,n}^2\geq j_{\ell+1/2,n}^2
 +p(p+1)-\ell(\ell+1).
 \label{cm:eq-angular-propagation}
\end{equation}
These are fixed-channel comparisons, not an invocation of unlabelled
global domain monotonicity.

We use the exact elementary enclosure
\begin{equation}
 \frac{333}{106}<\pi<\frac{22}{7}.
 \label{cm:eq-pi}
\end{equation}
For the lower bound, Machin's identity
\(\pi=16\arctan(1/5)-4\arctan(1/239)\) and alternating sums give a
strictly stronger rational bound; for the upper bound,
\[
 0<\int_0^1\frac{x^4(1-x)^4}{1+x^2}\,dx=\frac{22}{7}-\pi.
\]

For clarity we now derive every specialized Bessel-zero inequality used
below, apart from the one published first-zero estimate that is stated
explicitly before specialization.  The spherical Bessel identity and
recurrences are standard \cite{DLMF}.  Set
\[
 \mathsf j_\ell(x)=\sqrt{\frac{\pi}{2x}}J_{\ell+1/2}(x),
 \qquad
 P_\ell(x)=x^{\ell+1}\mathsf j_\ell(x).
\]
Then
\begin{equation}
 P_{\ell+1}=(2\ell+1)P_\ell-x^2P_{\ell-1},
 \qquad P_\ell'=xP_{\ell-1},
 \label{cm:eq-P-recurrence}
\end{equation}
with \(P_0=\sin x\) and \(P_1=\sin x-x\cos x\).  In particular,
\[
\begin{aligned}
P_2={}&(3-x^2)\sin x-3x\cos x,\\
P_3={}&(15-6x^2)\sin x+(x^3-15x)\cos x,\\
P_5={}&(945-420x^2+15x^4)\sin x
       +(-945x+105x^3-x^5)\cos x,\\
P_6={}&(10395-4725x^2+210x^4-x^6)\sin x\\
 &+(-10395x+1260x^3-21x^5)\cos x,\\
P_7={}&(135135-62370x^2+3150x^4-28x^6)\sin x\\
 &+(-135135x+17325x^3-378x^5+x^7)\cos x.
\end{aligned}
\]
Positive zeros are simple by ODE uniqueness.  At a zero of
\(\mathsf j_\ell\), the derivative recurrence gives
\(\mathsf j_{\ell+1}=-\mathsf j_\ell'\).  Combined with
\eqref{cm:eq-angular-propagation}, this gives strict interlacing and makes
the following sign/count ledger a complete enumeration, rather than a
sampled sign test:
\begin{center}
\small
\begin{tabular}{c|c|c|c}
endpoint & half-period & signs \(\mathsf j_0,\ldots,\mathsf j_\ell\)
& zero counts \(N_0,\ldots,N_\ell\)\\ \hline
\(77/10\)&\((2\pi,5\pi/2)\)&\(+,-\)&\(2,1\)\\
\(9\)&\((5\pi/2,3\pi)\)&\(+,+,-\)&\(2,2,1\)\\
\(10\)&\((3\pi,7\pi/2)\)&\(-,+,+,-,-,-,+\)&\(3,2,2,1,1,1,0\)\\
\(103/10\)&\((3\pi,7\pi/2)\)&\(-,+,+,-\)&\(3,2,2,1\)\\
\(21/2\)&\((3\pi,7\pi/2)\)&\(-,+\)&\(3,2\)\\
\(23/2\)&\((7\pi/2,4\pi)\)&\(-,-,+,+,-,-,-,+\)&\(3,3,2,2,1,1,1,0\)\\
\(61/5\)&\((7\pi/2,4\pi)\)&\(-,-,+\)&\(3,3,2\)\\
\(129/10\)&\((4\pi,9\pi/2)\)&\(+,-,-,+,+,-\)&\(4,3,3,2,2,1\)
\end{tabular}
\end{center}
For example, the endpoint signs follow from \eqref{cm:eq-pi}, the displayed
polynomials, and
\(\tan s>s\), \(\tan s>s+s^3/3\),
\(\tan s<s/(1-s^2/2)\) for \(0<s<1\).  The nonautomatic strict
comparisons can be taken as
\[
\begin{array}{c|c}
77/10&5\pi/2-x>3/20>10/77\\
21/2&7\pi/2-x>49/100>2/21\\
9&3\pi-x>21/50>9/26\\
61/5&4\pi-x>9/25>915/3646\\
103/10&7\pi/2-x>69/100>621540/938227\\
129/10&0<x-4\pi<17/50,\quad\tan(x-4\pi)<1700/4711\\
10&4/7<x-3\pi<29/50,\quad\tan(x-3\pi)<2900/4159\\
23/2&53/50<4\pi-x<\pi/2,\quad
\tan(4\pi-x)>546377/375000.
\end{array}
\]
All comparisons are between positive rationals and are checked by cross
multiplication.  The ledger proves
\begin{equation}
\begin{gathered}
 j_{3/2,2}>77/10,\quad j_{5/2,2}>9,\quad
 j_{7/2,2}>103/10,\quad j_{11/2,2}>129/10,\\
 j_{3/2,3}>21/2,\quad j_{5/2,3}>61/5,\\
 j_{13/2,1}>10,\quad j_{15/2,1}>23/2.
\end{gathered}
 \label{cm:eq-internal-zero-registry}
\end{equation}
This includes the requested higher-radial derivations.  For additional
detail, the second root of \(P_3\) is the unique root in
\((3\pi,7\pi/2)\): the derivative chain
\(R'=x\sin x\), \(P_2'=xR\), \(P_3'=xP_2\) first excludes every
earlier cell, and at \(x=103/10\), writing \(y=x-3\pi\), alternating
Taylor bounds give \(\tan y<5/4\), while
\[
 \frac{x(x^2-15)}{6x^2-15}-\frac32=\frac{5917}{621540}>0,
\]
so \(P_3(x)<0\).  Likewise, the third root of \(P_2\) is unique in
\((7\pi/2,4\pi)\); at \(x=61/5\),
\[
 4\pi-x>\frac{97}{265}>
 \frac{3x}{x^2-3},
\]
so \(P_2(x)>0\) before its decreasing zero.  At \(x=129/10\), the
explicit \(P_5\) coefficients and
\(\tan(x-4\pi)<3/8<|B|/A\) give \(P_5(x)<0\); angular propagation of
\(j_{5/2,3}>61/5\) puts the third \(P_5\)-zero above this point, hence
its second zero is also above it.  These arguments use no root table.

For the remaining first zeros we use Lorch's strict first-zero inequality
\cite{Lorch1993}
\[
 j_{\nu,1}^2>
 \frac{24(\nu+1)^2}
 {1-2\nu+\sqrt{(2\nu+3)(2\nu+11)}}-2(\nu^2-1),
 \qquad \nu>-1.
\]
At \(\nu=\ell+1/2\), its right side is
\[
 R_\ell=\frac{3(2\ell+3)\sqrt{(\ell+2)(\ell+6)}
 -2\ell^2+\ell+6}{4}.
\]
Exact positive-side squaring gives
\begin{center}
\small
\begin{tabular}{c|c|c}
\(\ell\)&first-zero wall&integer residual proving \(R_\ell>c^2\)\\ \hline
2&\(51/10\)&\(105/4-(51/10)^2=6/25\)\\
3&\(13/2\)&\(81^2\cdot5-178^2=1121\)\\
4&\(15/2\)&\(66^2\cdot15-247^2=4331\)\\
5&\(87/10\)&\(198189\)\\
8&\(71/6\)&\(1026^2\cdot35-6067^2=35171\)\\
9&\(64/5\)&\(1575^2\cdot165-20059^2=6939644\)\\
10&\(69/5\)&\(3450^2\cdot3-5911^2=767579\).
\end{tabular}
\end{center}
Together with \eqref{cm:eq-internal-zero-registry}, this yields the complete
first-mode walls
\begin{equation}
\begin{array}{c|cccccccccc}
\ell&1&2&3&4&5&6&7&8&9&10\\ \hline
\beta_{\ell,1}&4&51/10&13/2&15/2&87/10&10&23/2&71/6&64/5&69/5.
\end{array}
 \label{cm:eq-first-zero-table}
\end{equation}
The \(\ell=1\) wall follows internally from \(\tan x=x\): the first
positive root lies in \((\pi,3\pi/2)\), and \(\tan4<2<4\).
Applying \eqref{cm:eq-angular-propagation} to the higher-radial seeds gives,
for later use,
\begin{align}
 j_{p+1/2,2}^2&>9409/100+p(p+1) &&(p\geq3),\label{cm:eq-prop-second-a}\\
 j_{p+1/2,3}^2&>3571/25+p(p+1) &&(p\geq2),\label{cm:eq-prop-third}\\
 j_{p+1/2,2}^2&>13641/100+p(p+1) &&(p\geq5).
 \label{cm:eq-prop-second-b}
\end{align}
Equations \eqref{cm:eq-shell-ball}--\eqref{cm:eq-prop-second-b} transfer
every displayed wall to the shell at the same radial index.

\subsection{The first and next angular staircases}

Throughout the staircase and optical calculations below we use the
abbreviation
\[
 z:=z_\rho=\frac{\pi}{1-\rho}.
\]

\begin{lemma}[First two angular bands]
\label{cm:lem-first-bands}
For \(\rho_c\leq\rho\leq7/8\) and
\(z_\rho\leq K\leq k_2(\rho)\), one has
\(N_D(A_{\rho,1},K^2)<W(\rho,K)\).
\end{lemma}

\begin{proof}
Since \(z_\rho>3\), \eqref{cm:eq-product-bound} excludes all \(n\geq2\)
and all \(\ell\geq2\) through \(k_2\).  Thus the strict caps are
\[
 N_D=0\quad(K=z_\rho),\qquad
 N_D\leq1\quad(z_\rho<K\leq k_1),
 \qquad N_D\leq4\quad(k_1<K\leq k_2).
\]
At \(K=k_m\), the newly possible channel \(m\) is excluded because the
count is strict.  At the first wall,
\[
 W(\rho,z_\rho)=\frac{2\pi^2}{9}
 \frac{1+\rho+\rho^2}{(1-\rho)^2}>\frac{2\pi^2}{9}>2>1.
\]
For \(F(\rho)=W(\rho,k_1(\rho))\), direct logarithmic differentiation gives
\[
 \frac{F'}F=\frac{1+2\rho}{1+\rho+\rho^2}
 +\frac{2(\pi^2-(1-\rho)^2)}{(1-\rho)(\pi^2+2(1-\rho)^2)}>0.
\]
At \(\rho_c\), exact substitution of
\(157/50<\pi<22/7\) gives
\[
 F(\rho_c)>
 \frac{388430104857}{92514296875}
 =4+\frac{18372917357}{92514296875}>4.
\]
Strict increase in \(K\) pays both caps and proves the lemma.
\end{proof}

The closed band of Lemma~\ref{cm:lem-first-bands} meets
\(\mathcal D_{16}\) exactly in
\(\rho_c\leq\rho<7/8\), \(z_\rho<K\leq k_2\).  Indeed
\(k_2<K_0\) follows from \(K_0>64\) and \(k_2<27\), and on the seam
\(k_2<27<1944\).  Thus exact subtraction gives
\begin{equation}
\begin{aligned}
 \mathcal D_{17}={}&
 \{\rho_*<\rho<\rho_c:(2\rho)^{-1}<K<U(\rho)\}\\
 &\cup\{\rho_c\leq\rho<7/8:k_2(\rho)<K<U(\rho)\}.
\end{aligned}
 \label{cm:eq-D17}
\end{equation}
The two former pilot rectangles
\[
\begin{aligned}
 B_0&=[999/2000,1001/2000]\times[67/10,168/25],\\
 B_1&=[999/2000,1001/2000]\times[168/25,673/100]
\end{aligned}
\]
are analytically subsumed.  Monotonicity of \(z_\rho,k_2\) and
\eqref{cm:eq-pi} give
\[
 \max z_\rho<\frac{44000}{6993}<\frac{67}{10},
\]
\[
 \min k_2^2-(168/25)^2>\frac{126027526}{626250625}>0,
 \quad
 \min k_2^2-(673/100)^2>\frac{668749071}{10020010000}>0.
\]
Hence no pilot box is subtracted a second time.

\begin{lemma}[Staircase through \(k_5\)]
\label{cm:lem-k5}
For \(\rho_c\leq\rho\leq7/8\) and
\(k_2(\rho)<K\leq k_5(\rho)\), one has \(N_D<W\).
\end{lemma}

\begin{proof}
Put
\[
 c_2=51/10,\qquad c_3=13/2,\qquad c_4=15/2,
 \qquad q_m(\rho)=\max\{k_m(\rho),c_m\}.
\]
The first-zero walls in \eqref{cm:eq-first-zero-table} and
\eqref{cm:eq-product-bound} give the successive caps
\[
 4,\quad 9,\quad 16,\quad 25
\]
on \(K\leq q_2\), \(q_2<K\leq q_3\),
\(q_3<K\leq q_4\), and \(q_4<K\leq k_5\).  Since
\(z\geq\pi+1/2>7/2\), \(4z^2>z^2+30\), so no second radial mode is
present.  The cumulative multiplicities are
\(1,1+3,1+3+5,1+3+5+7,1+3+5+7+9\).

For \(F_a(\rho)=(1-\rho^3)(z_\rho^2+a)^{3/2}\), the sign of
\(F_a'\) is the sign of
\[
 z_\rho^2(1+\rho)-a\rho^2.
\]
It is positive for \(a=6,12,20\) on \([\rho_c,7/8]\).  Thus moving
payments increase with \(\rho\), while fixed-frequency payments decrease.
The following exact split payments cover both one-sided traces:
\begin{center}
\small
\begin{tabular}{c|c|c|c}
cap&split&fixed wall on the left&moving wall on the right\\ \hline
9&\(3/10\)&\(W(3/10,51/10)>9\)&\(W(3/10,k_2)>9\)\\
16&\(1/2\)&\(W(1/2,13/2)=107653/6336>16\)&\(W(1/2,k_3)>196/11>16\)\\
25&\(1/2\)&\(W(1/2,15/2)=165375/6336>25\)&\(W(1/2,k_4)>101528/3960>25\).
\end{tabular}
\end{center}
For cap \(9\), the exact fixed and moving cross-product reserves are
\(12485961\) and \(592618642\), respectively.  Cap \(4\) is paid by
\[
 W(\rho_c,k_2)>\frac7{99}\frac{342}{343}\,73>4.
\]
All inequalities are strict after a channel enters; equality at a zero wall
retains the lower cap.

It remains to compare with the upper owner.  If \(\rho\leq1/2\), then
\(\eta_\rho=\omega_0<1/9\): indeed \(9\sqrt3<5\pi\) follows after
positive squaring from \(\pi>333/106\) and
\(25\cdot333^2-243\cdot106^2=41877>0\).  Since \(k_5\) increases with
\(\rho\),
\[
 k_5(\rho)^2\leq k_5(1/2)^2=4\pi^2+30
 <4\left(\frac{22}{7}\right)^2+30
 =\frac{3406}{49}<81,
\]
and therefore \(\eta_\rho k_5(\rho)<1\).

If \(\rho\geq1/2\), write \(\rho=\cos\theta\).  Then
\(\pi\eta_\rho=\sin\theta-\theta\cos\theta
\leq\theta(1-\rho)\), while \(\theta\leq\pi/3\) and
\(\sqrt{z^2+30}<z+15/z\).  Hence
\[
 \eta_\rho k_5(\rho)
 <\frac\pi3+\frac{5(1-\rho)^2}{\pi}
 \leq\frac\pi3+\frac5{4\pi}
 <\frac{41}{28}<\frac{131}{75}<C_0.
\]
The last chain uses \(3<\pi<22/7\) and
\(\sqrt2>7/5\), the latter because \(2>49/25\).  Thus in both ratio cases
\(\eta_\rho k_5<C_0\).
But \(\eta_\rho K_0=C_0+\sqrt{a_\rho K_0}>C_0\), hence
\(k_5<K_0\).  On the seam, with \(t=1-\rho\leq1/6\),
\[
 t^4k_5^2=t^2\pi^2+30t^4<65/216<2916=t^4U^2.
\]
This proves the lemma on all faces.
\end{proof}

Consequently
\begin{equation}
\begin{aligned}
 \mathcal D_{18}={}&
 \{\rho_*<\rho<\rho_c:(2\rho)^{-1}<K<U(\rho)\}\\
 &\cup\{\rho_c\leq\rho<7/8:k_5(\rho)<K<U(\rho)\}.
\end{aligned}
 \label{cm:eq-D18}
\end{equation}

\subsection{The two-sided staircase}

Put
\[
 \rho_0=\frac1{\sqrt{337}},\qquad d=\frac{\sqrt{337}}2,
 \qquad p_1(\rho)=\sqrt{4z_\rho^2+2}.
\]

\begin{lemma}[High side through \(k_6\)]
\label{cm:lem-k6}
For \(\rho_c\leq\rho\leq7/8\) and
\(z_\rho\leq K\leq k_6(\rho)\), one has \(N_D<W\).
\end{lemma}

\begin{proof}
For \(\rho\leq1/4\), the only possible new modes beyond the first modes
\(\ell=0,\ldots,4\) are \((0,2),(1,2)\); the first \(\ell=5\) mode is
delayed by \(j_{11/2,1}>17/2\).  For \(\rho\geq1/4\),
\(z>4\) makes every second mode absent through \(k_6\), and first modes
\(\ell=0,\ldots,5\) are the complete inventory.  The cumulative first-mode
caps are \(1,4,9,16,25,36\); the two possible second modes give caps
\(26,29\).  Their entry walls are
\[
\begin{array}{c|c}
(1,1)&k_1\\
(2,1)&\max\{k_2,51/10\}\\
(3,1)&\max\{k_3,13/2\}\\
(4,1)&\max\{k_4,15/2\}\\
(5,1)&\max\{k_5,17/2\}\\
(0,2)&2z\\
(1,2)&\max\{p_1,77/10\}.
\end{array}
\]
For a moving wall \(T^2=az^2+b\), \(W(\rho,T)\) increases whenever
\(a\pi^2(1+\rho)>b\rho^2(1-\rho)^2\), which holds here because
\(b/a\leq42\).  Fixed walls decrease with \(\rho\).  Exact payments are
obtained by the splits
\[
 (k_1,1/8,4),\ (k_2,3/10,9),\ (k_3,1/2,16),\
 (k_4,1/2,25),\ (k_5,13/25,36),\ (p_1,1/5,29).
\]
Writing
\(A=a(333/106)^2/(1-r)^2+b\) and
\(B=7(1-r^3)/99\), each moving payment is the positive rational check
\(B^2A^3-C^2>0\).  In the displayed order the numerators are
\[
\begin{gathered}
153864824754340538785,
1399810848073457853053,
137027505353736631,\\
2507119282010251109,
92672404686738742434741,
373255772037478993002833,
\end{gathered}
\]
over positive denominators.  The corresponding fixed-side margins are
positive as well; for example they are
\(1387329/11000000,6277/6336,775/704,452843/343750\), and
\(849901/281250\) at the nontrivial splits.  Finally
\(W(\rho_c,2z)>26\).  Every possible mode is therefore paid, including
coincident entries, and every equality remains assigned to the lower cap.
 Parts~I, II, and IV of the \emph{Purely Analytic Supplement} map every
 printed moving numerator and fixed margin to its wall and cap.
\end{proof}

\begin{lemma}[Lower side through \(d\)]
\label{cm:lem-lower-d}
For \(\rho_0<\rho<\rho_c\) and
\((2\rho)^{-1}<K\leq d\), one has \(N_D<W\).
\end{lemma}

\begin{proof}
Because \(3z>d\), every \(n\geq3\) mode is absent.  The zero registry and
angular propagation exclude first modes \(\ell\geq6\) and second modes
\(\ell\geq3\).  The complete cap table is
\[
\begin{array}{c|c}
\text{open-lower, closed-upper range}&N_D\text{ cap}\\ \hline
L<K\leq4&1\\
4<K\leq51/10&4\\
51/10<K\leq13/2,\ K\leq2z&9\\
51/10<K\leq13/2,\ K>2z&10\\
13/2<K\leq15/2,\ K\leq2z&16\\
13/2<K\leq15/2,\ K>2z&17\\
15/2<K\leq77/10&26\\
77/10<K\leq17/2&29\\
17/2<K\leq9&40\\
9<K\leq d&45.
\end{array}
\]
In fact \(13/2<2z<15/2\), so the cap-10 row is empty.  The fixed-wall
lower bounds, using \(\rho<\rho_c<7/50\), exceed their caps by
\[
 \frac{793292}{1546875},\frac{486236661}{1375000000},
 \frac{333100003}{99000000},\frac{329797}{88000},
 \frac{3590476297}{1125000000},\frac{327078887}{99000000},
 \frac{8805519}{1375000},
\]
respectively.  The initial cap is paid by
\(W(\rho_c,L(\rho_c))>19/6>1\), and the moving branch by
\(W(1/19,2z)>508/27>17\).  This proves every row, including \(K=d\).
 Parts~I and IV of the \emph{Purely Analytic Supplement} record each margin
 with its corresponding row and prove the moving \(2z\) payment.
\end{proof}

The upper-floor comparisons are strict.  On the lower side,
\(H_0(\rho_*)=(2\rho_*)^{-1}>d\), while
\(K_0>C_0/\omega_0>12>d\).  On the high side, \(K_0>12>k_6\) for
\(\rho\leq1/2\), \(K_0>384>k_6\) for \(1/2<\rho<5/6\), and the seam
owner exceeds \(54>k_6\).  Thus exact subtraction gives
\begin{equation}
\boxed{\begin{aligned}
\mathcal D_{19}={}&
\{\rho_*<\rho\leq\rho_0:(2\rho)^{-1}<K<U(\rho)\}\\
&\cup\{\rho_0<\rho<\rho_c:d<K<U(\rho)\}\\
&\cup\{\rho_c\leq\rho<7/8:k_6(\rho)<K<U(\rho)\}.
\end{aligned}}
\label{cm:eq-D19}
\end{equation}
At \(\rho=\rho_0\), \((2\rho_0)^{-1}=d\), so the would-be new lower
fiber is empty and the first component owns the equality slice.

\subsection{Combined closure of the lower components and the high staircase}

Set
\[
 \sigma:=\frac{3\omega_0}{4}.
\]
We make every comparison in the required ordering explicit.  The
alternating Machin sums give
\[
\begin{aligned}
 \pi
 &<16\left(\frac15-\frac1{3\cdot5^3}
 +\frac1{5\cdot5^5}-\frac1{7\cdot5^7}
 +\frac1{9\cdot5^9}\right)\\
 &\quad-4\left(\frac1{239}-\frac1{3\cdot239^3}\right)
 <\frac{355}{113},
\end{aligned}
\]
and the last gap is exactly
\[
 \frac{45167474711}{189820334689453125}>0.
\]
Since
\[
 3\cdot1101^2\cdot113^2-607^2\cdot355^2=1998482>0,
\]
we have
\[
 \sqrt3>\frac{607\cdot355}{1101\cdot113}
 >\frac{607\pi}{1101},
\]
and therefore
\begin{equation}
 \omega_0=\frac{\sqrt3}{2\pi}-\frac16>\frac{40}{367}.
 \label{cm:eq-omega-lower}
\end{equation}
On the other side,
\[
 25\cdot333^2-243\cdot106^2=41877>0
\]
and \(\pi>333/106\) imply \(9\sqrt3<5\pi\), hence
\(\omega_0<1/9\).  Moreover
\(2\cdot32^2-45^2=23>0\), so
\(\sqrt2>45/32\), \(C_0>7/4\), and
\[
 \rho_*=\frac{\omega_0}{2C_0}<\frac2{63}.
\]
Now all remaining comparisons are witnessed by
\[
\begin{gathered}
 4\cdot337<63^2,\qquad 18^2<337,\qquad
 \frac1{18}<\frac3{40}<\frac{3\omega_0}{4}
 <\frac1{12}<\frac7{51},\\
 \frac7{51}<\frac1{1+2\pi}<\frac17,\qquad
 \frac17<\frac{39}{50}<\frac78.
\end{gathered}
\]
Here the first inequality on the second line uses \(\pi<22/7\), and the
second uses \(\pi>3\).  Thus, without decimal comparison,
\[
 \rho_*<\rho_0<\sigma<\rho_c<\frac{39}{50}<\frac78.
\]

\begin{lemma}[Closure of both lower components of \(\mathcal D_{19}\)]
\label{cm:lem-lower-closure}
The strict inequality \(N_D<W\) holds on the first two components of
\eqref{cm:eq-D19}.
\end{lemma}

\begin{proof}
We divide the proof into a wedge, a finite staircase, and one remaining
floor cell.

For a low-interface tail beginning at \(x_r=r+1/2<\mu=\rho K\), put
\[
 n_r=\lfloor\mu-x_r\rfloor,\qquad
 q_r=x_r+n_r,\qquad p=\lfloor\omega_0K\rfloor.
\]
The concave-head/convex-tail split used in
\eqref{cm:eq-seam-tail}, now with its coarse first-shelf payment, gives
\begin{equation}
 \frac{\mathcal T_r}{2}
 <\int_{x_r}^K A_{\rho,K}(z)\,dz
 +\delta_r-\frac{p-n_r}{4},
 \qquad
 0\leq\delta_r<
 \frac{2(\mu-q_r)^{5/2}}{15\sqrt\mu}.
 \label{cm:eq-lower-wedge-tail}
\end{equation}
When \(n_r=0\), the concave head is absent and the formula follows from
the convex tail alone.  From \eqref{cm:eq-omega-lower} and
\(400\cdot337-367^2=111>0\),
\[
 \omega_0d>
 \frac{40}{367}\frac{\sqrt{337}}2
 =\frac{20\sqrt{337}}{367}>1.
\]
If
\(0<\rho\leq\sigma\), \(K>d\), and \(\rho K>1/2\), then with
\(x=\omega_0K>1\) and \(m=\lfloor x\rfloor\),
\[
 n_r\leq\left\lfloor\frac{3x}{4}-\frac12\right\rfloor<m=p.
\]
Since \(\delta_r<2\sqrt2/15<1/4\), every tail in
\eqref{cm:eq-lower-wedge-tail} is strict.  The hypotheses hold throughout
the first component of \(\mathcal D_{19}\) and on
\(\rho_0\leq\rho\leq\sigma\), \(K>d\).  This includes the complete
\(\rho=\rho_0\) fiber.

For \(\rho_0<\rho<\rho_c\) and \(K\leq\sqrt{114}\), the complete zero
registry and \eqref{cm:eq-angular-propagation} give the following exhaustive
inventory:
\begin{center}
\small
\begin{tabular}{c|l|c}
range (lower face open, upper face closed)&newly possible modes&cap\\ \hline
\(d<K\leq10,\ K\leq3z\)&first \(0{:}5\), second \(0{:}2\)&45\\
\(d<K\leq10,\ K>3z\)&preceding plus \((0,3)\)&46\\
\(10<K\leq81/8\)&plus \((6,1)\), possibly \((0,3)\)&59\\
\(81/8<K\leq21/2\)&plus \((3,2)\)&66\\
\(21/2<K\leq\sqrt{7073}/8\)&plus \((1,3)\)&69\\
\(\sqrt{7073}/8<K\leq\sqrt{114}\)&plus \((4,2)\)&78.
\end{tabular}
\end{center}
First modes \(\ell\geq7\), second modes \(\ell\geq5\), third modes
\(\ell\geq2\), and all \(n\geq4\) are absent.  In particular,
\[
 (81/8)^2+8=7073/64
\]
is the propagated \((4,2)\) wall.  Since \(\rho<\rho_c<1/7\), the
uniform estimate
\[
 W(\rho,K)>
 \frac{2}{9(22/7)}\left(1-\frac1{7^3}\right)K^3
\]
pays the successive caps with exact positive reserves
\[
 \frac{2908}{539},\quad\frac{6199}{539},\quad
 \frac{990435}{137984},\quad\frac{555}{44},\quad\frac{159}{44}.
\]
The first number pays even cap \(46\) from \(K>9\); the last pays cap
\(78\) from \(K>21/2\).  Hence \(K=\sqrt{114}\) is included.

It remains to consider \(\sigma<\rho<\rho_c\) and
\(K>\sqrt{114}\).  Put
\[
 p=\lfloor\omega_0K\rfloor,\qquad
 m=\lfloor\rho K-1/2\rfloor.
\]
Then \(p\geq1\).  Except on
\begin{equation}
 \mathcal B=\left\{\sigma<\rho<\rho_c,\quad K>\sqrt{114},\quad
 \rho K\geq5/2,\quad K<2/\omega_0\right\},
 \label{cm:eq-exceptional-B}
\end{equation}
one has \(m\leq2p-1\).  Indeed this is immediate when \(\rho K<5/2\);
when \(K\geq2/\omega_0\), one has \(p\geq2\).  Moreover,
\[
 \rho_c<\frac17<\frac{60}{367}<\frac{3\omega_0}{2},
\]
where the first and last comparisons were proved above and
\(367<420\).  Hence \(\rho K-1/2<2p\).

Away from a half-integer \(\mu\), all interface remainders in
\eqref{cm:eq-lower-wedge-tail} are equal.  Summing over
\(r=0,\ldots,m\) gives the average floor reserve
\[
 \frac14\left(p-\frac m2\right)\geq\frac18.
\]
For \(m=0\), use \(\delta<1/4\); for \(m\geq1\), \(\mu>3/2\) and
\(\delta<2/(15\sqrt{3/2})<1/8\).  At
\(\mu\in\mathbb Z+1/2\), assign the interface to the high side; then the
remainder is zero and strictness remains in the nonzero convex tail.

On \(\mathcal B\), \(\omega_0>40/367\), so \(K<367/20=:R\).  The
turning estimate
\[
 A_{\rho,K}(z)\leq G_K(z)<\frac{(K-z)^{3/2}}{3\sqrt K}
\]
and monotonicity in \(K\) reduce the proxy count at
\(z_\ell=\ell+1/2\) to
\[
 \frac{(357-20\ell)^3}{1321200}
 <\left(c_\ell+\frac34\right)^2.
\]
The integer numerators of these strict comparisons are
\begin{center}
\small
\begin{tabular}{c|rrrrrrrrrrrrrrr}
\(\ell\)&0&1&2&3&4&5&6&7&8&9&10&11&12&13&14\\ \hline
\(c_\ell\)&6&5&5&4&4&3&3&3&2&2&1&1&1&1&0
\end{tabular}
\end{center}
and, in the same order,
\[
\begin{gathered}
14697882,5409422,11827162,3611502,8555642,
1604782,5267322,8361062,\\
2346202,4446342,176282,1474822,2444562,3133502,286642.
\end{gathered}
\]
The \(\ell=14\) row and monotonicity remove the entire angular tail.
Therefore the full multiplicity cap is
\[
 \sum_{\ell=0}^{13}(2\ell+1)c_\ell=395.
\]
But \(\rho K\geq5/2\), \(\rho<\rho_c<11/80\), and
\(1/\pi>113/355\) give
\[
 W(\rho,K)>
 \frac{160293325}{378004}
 =395+\frac{10981745}{378004}>395.
\]
Thus \(\mathcal B\) is closed as well.  The face \(\rho K=5/2\) is
\(\mathcal B\)-owned and \(K=2/\omega_0\) is aggregate-owned.  All lower
components of \(\mathcal D_{19}\) are proved.
\end{proof}

We now turn to the high component.  Define the strict channel proxy
\begin{equation}
 b_{\ell,n}(z)^2=
 \max\{n^2z^2+\ell(\ell+1),\,\beta_{\ell,n}^2\},
 \label{cm:eq-channel-proxy}
\end{equation}
For the radial channel we set, for every \(n\geq1\),
\[
 \beta_{0,n}:=j_{1/2,n}=n\pi,
\]
the last identity following from
\(J_{1/2}(x)=\sqrt{2/(\pi x)}\sin x\).  For \(\ell\geq1\), the
first-mode \(\beta\)'s are in
\eqref{cm:eq-first-zero-table} and
\[
\begin{aligned}
 \beta_{1,2}&=77/10,&
 \beta_{2,2}&=9,&
 \beta_{3,2}&=103/10,&
 \beta_{4,2}&=\sqrt{11409}/10,\\
 \beta_{5,2}&=129/10,&
 \beta_{1,3}&=21/2,&
 \beta_{2,3}&=61/5.
\end{aligned}
\]
By \eqref{cm:eq-product-bound} and \eqref{cm:eq-shell-ball}, a mode can
be counted only if \(b_{\ell,n}(z)<K\).

\begin{lemma}[Closed high staircase through \(k_{11}\)]
\label{cm:lem-k11}
For \(\rho_c\leq\rho\leq7/8\) and
\(z_\rho\leq K\leq k_{11}(\rho)\), one has \(N_D<W\).
\end{lemma}

\begin{proof}
The exhaustive checkpoint inventories are
\begin{center}
\small
\begin{tabular}{c|l}
checkpoint&possible modes\\ \hline
\(k_7\)&first \(0{:}6\), second \(0{:}1\), no \(n\geq3\)\\
\(k_8\)&first \(0{:}7\), second \(0{:}3\), no \(n\geq3\)\\
\(k_9\)&first \(0{:}8\), second \(0{:}4\), no \(n\geq3\)\\
\(k_{10}\)&first \(0{:}9\), second \(0{:}5\), third \(0{:}1\), no \(n\geq4\)\\
\(k_{11}\)&first \(0{:}10\), second \(0{:}4\), third \(0{:}1\), no \(n\geq4\).
\end{tabular}
\end{center}
For example, at \(k_{11}\) the mode \((5,2)\) would require both
\(z^2<34\) from the product wall and
\(z^2>(129/10)^2-132=3441/100\) from the ball wall, a contradiction.
Also \(4z>k_{11}\) at \(z\geq\pi+1/2\).  Thus radial indices
\(1,2,3,4\) and angular indices through \(11\) are proved above to be
exhaustive for the stated checkpoints.

Writing \(H\) for the highest first channel and \(h\) for the highest
second channel, the block cap is \((H+1)^2+(h+1)^2\), with an additional
\(1\) or \(4\) for third modes.  The complete cap tables at \(k_8\) and
\(k_9\) are
\[
\begin{array}{c|ccccc}
k_8:H\backslash h&\varnothing&0&1&2&3\\ \hline
7&64&65&68&73&80\\
6&49&50&53&-&-\\
5&36&37&40&45&52\\
\leq4&25&26&29&34&41
\end{array}
\qquad
\begin{array}{c|cccccc}
k_9:H\backslash h&\varnothing&0&1&2&3&4\\ \hline
8&81&-&-&-&-&-\\
7&64&65&68&73&-&-\\
6&49&50&53&58&65&74\\
5&36&37&40&45&52&61\\
\leq4&25&26&29&34&41&50.
\end{array}
\]
The dashes follow from incompatible necessary inequalities, not from a
numerical sample.  At \(k_{10}\) the cap lists are
\[
 100;\quad81,97;\quad64,80,89,100;\quad74,69,78,85,89,
\]
and at \(k_{11}\) they are
\[
 121;\quad100,125;\quad106,110;\quad93.
\]

Define the exact finite proxy
\[
 C_z(K)=\sum_{(\ell,n)}(2\ell+1)
 \mathbf1_{\{b_{\ell,n}(z)<K\}},
\]
where the sum runs over the finite inventory above.  Between proxy events
it is constant and \(W\) increases in \(K\).  Every squared event is either
constant or affine in \(z^2\).  At a moving event \(K^2=sz^2+d_0\), the
derivative of its Weyl payment has the sign of
\begin{equation}
 s\pi^2\frac{1+\rho}{(1-\rho)^2}-\rho^2d_0.
 \label{cm:eq-moving-payment-sign}
\end{equation}
Thus the finite split registry
\[
\frac15,\frac3{10},\frac13,\frac38,\frac25,\frac5{12},\frac37,
\frac12,\frac{107}{200},\frac8{15},\frac35,\frac{16}{25},
\frac{13}{20},\frac23,\frac4{25},\frac14,\frac7{20},
\]
together with \(z^2=16,68/3,34\), reduces every cell to exact rational
endpoint arithmetic.  The smallest representative reserves in the six
successive bands are
\[
 \frac{1387329}{11000000},\quad
 \frac{849901}{281250},\quad
 \frac{835355}{577368},\quad
 \frac{835355}{577368},\quad
 \frac{943655}{171072},\quad
 \frac{507845}{261954}.
\]
Every realizable event is strictly paid, with coincident events charged at
their combined multiplicity.

There is one conservative row requiring explicit treatment.  Angular
propagation gives
\[
 j_{9/2,2}^2>\left(\frac{103}{10}\right)^2+8
 =\frac{11409}{100}.
\]
If \((4,2)\) were counted at \(k_9\), this ball wall would force
\(z^2>2409/100\), whereas the product wall forces \(z^2<70/3\); since
\(2409/100-70/3=227/300>0\), the printed cap-74 cell is actually empty.
If one deliberately retains the weaker conservative cell, counting the
mode forces \(K>207/20\) and \(\rho<7/20\).  Therefore
\begin{equation}
 W(\rho,K)>
 \frac{52823261673}{704000000}
 =74+\frac{727261673}{704000000}>74.
 \label{cm:eq-cap74}
\end{equation}
This is the corrected cap-74 payment; it is smaller than the generic
\(k_9\) summary reserve and must not be hidden by a claim that all omitted
cells have larger reserves.

At every \(K=k_m\), the \((m,1)\) mode is on its product equality wall
and is excluded.  At every fixed Bessel wall the zero inequality is strict.
The event calculation therefore covers all equality and coincidence faces,
and proves the lemma.
\end{proof}

The combined finite analytic ledger---covering the zero/sign registry, the
two-sided and lower staircases, the lower closure, the high staircase,
optical constants, and exact subtraction---can be replayed without floating
point.  Its core contains \(488\) substantive rational/radical checks and
\(65\) bookkeeping identities; these are whole-ledger counts, not a count of
\(k_{11}\)-only cells.  It also checks the sign of
\eqref{cm:eq-moving-payment-sign} on the named faces.  Every proof decision
is a cross multiplication or positive-side squaring; decimal displays play
  no role.  The disjoint (611)-row allocation and every hand-checkable
  witness are printed in Parts~I--VI of the \emph{Purely Analytic Supplement}.

\subsection{The optical closure and the exact residual
\texorpdfstring{\(\mathcal D_{20}\)}{D20}}

\begin{lemma}[All-frequency optical theorem]
\label{cm:lem-optical}
For \(39/50\leq\rho<1\) and \(K\geq0\),
\(N_D(A_{\rho,1},K^2)\leq W(\rho,K)\), with strict inequality for
\(K>0\).
\end{lemma}

\begin{proof}
Put \(\varepsilon=1-\rho\leq11/50\),
\(a=\varepsilon K\), and \(c=1126/625\).  We use two inclusive screens
meeting at \(a=c/\varepsilon\).

For the product screen, \eqref{cm:eq-product-bound} implies that a counted
mode satisfies
\[
 n^2\pi^2+\varepsilon^2\ell(\ell+1)<a^2.
\]
There is no mode for \(a\leq\pi\).  For \(a>\pi\), write
\(a/\pi=N+t\), with \(N\geq1\), \(0<t\leq1\), assigning
\((N,t)=(m-1,1)\) at \(a=m\pi\).  If
\[
 M_n=\left\lceil\sqrt{(b_n/\varepsilon)^2+1/4}-1/2\right\rceil,
\]
the exact angular multiplicity is bounded by
\[
 \mathcal P_\varepsilon(a)=\sum_{n=1}^N M_n^2.
\]
Direct summation of the continuous-minus-discrete product deficit gives
\begin{equation}
 \frac{D(a)}{\pi^2}=\frac{N^2}{2}
 +N\left(t^2+\frac16\right)+\frac{2t^3}{3}.
 \label{cm:eq-optical-product-deficit}
\end{equation}
For \(C_D=1382/3125\), the difference
\(D/\pi^2-C_D(N+t)^2\) increases in \(N\), and at \(N=1\) is minimized
at \(t=C_D\), where it equals
\(1822532/91552734375>0\).  Thus \(D>C_Da^2\).

The elementary ceiling bound and quarter-circle integral give
\[
 \varepsilon^2\mathcal P_\varepsilon
 <S_0+\varepsilon a^2/4+\varepsilon^2a/\pi.
\]
On \(a\leq c/\varepsilon\), using \(q=106/333>1/\pi\) and
\(\varepsilon\leq11/50\), the exact reserve is
\begin{equation}
 C_D-\left(\frac{2cq}{3}+\frac{11}{200}
 +(11/50)^2q^2\right)
 =\frac{39569}{2772225000}>0.
 \label{cm:eq-optical-low-reserve}
\end{equation}
This proves the low screen, including all radial and angular equality walls.

For the complementary-action screen, use the inverse-action construction
proved in Lemma~\ref{thin:lem-action-geometry}: let the decreasing shell
action have inverse \(R_{\mathcal A}\), put \(F=R_{\mathcal A}^2\), and
\(g=-F'\).  If \(\tau\) is the actual curvature interface and
\(T=a/\pi\), that lemma gives
\[
 g\text{ decreasing on }(0,\tau),\quad
 g\text{ increasing on }(\tau,T),
\]
\[
 F(0)=a^2,\quad F(\tau)=\rho^2a^2,\quad F(T)=0,\quad
 g(T)=2\pi\rho a.
\]
For \(C(t)=\#\{n:n-1/4<t\}\), set \(w=t-C(t)\) and
\(\Psi(t)=\int_0^tw(s)\,ds-t/4\).  The complete sawtooth and Stieltjes
calculation in Lemma~\ref{thin:lem-radial-deficit}, including the ungridded
curvature interface \(\tau\), gives
\begin{equation}
 D_{\rm rad}:=\int_0^TF(t)\,dt
 -\sum_{n-1/4<T}F(n-1/4)
 \geq\frac{\rho^2a^2}{4}-\frac{\pi\rho a}{4}.
 \label{cm:eq-optical-radial-deficit}
\end{equation}
The exact half-integer angular count from
Lemma~\ref{thin:lem-layer-cake} satisfies
\(\varepsilon^2M_\varepsilon(x)^2<(x+\varepsilon/2)^2\), even on an
angular wall.  Since fewer than \(a/\pi+1/4\) radial layers occur, the
angular error is strictly below
\begin{equation}
 E=\left(\frac a\pi+\frac14\right)
 \left(\varepsilon a+\frac{\varepsilon^2}{4}\right).
 \label{cm:eq-optical-angular-error}
\end{equation}
On \(a\geq c/\varepsilon\), the lower screen from
\eqref{cm:eq-optical-radial-deficit} decreases and the upper screen from
\eqref{cm:eq-optical-angular-error} increases with \(\varepsilon\).  At
\(\varepsilon=11/50\), their exact difference is
\begin{equation}
 \frac{14817541}{472867032960000}>0.
 \label{cm:eq-optical-high-reserve}
\end{equation}
Thus the strict layer-cake sum is below the action integral, and the phase
bridge proves the high screen.  Both screens include
\(a=c/\varepsilon\); at \(K=0\) both sides vanish.
\end{proof}

Combining Lemmas~\ref{cm:lem-lower-closure}, \ref{cm:lem-k11}, and
\ref{cm:lem-optical} with \eqref{cm:eq-D19} leaves exactly
\begin{equation}
 \boxed{
 \mathcal D_{20}=
 \left\{(\rho,K):\rho_c\leq\rho<\frac{39}{50},\quad
 k_{11}(\rho)<K<U(\rho)=K_0(\rho)\right\}.}
 \label{cm:eq-D20}
\end{equation}
Indeed \(k_{11}<U\): on \(\rho\geq\rho_c\),
\(a_\rho\geq1\) and \(\eta_\rho<1/8\), so \(K_0>64\), whereas
\(k_{11}<64\) for \(\rho<7/8\); the seam owner is greater than \(1944\).
On the surviving ratio range neither \(H_0\) nor the seam branch is
eligible, so \(U=K_0\).  The face \(K=k_{11}\) belongs to the closed
staircase, \(\rho=39/50\) belongs to the optical theorem, and \(K=U\)
belongs to the fixed-ratio high-frequency owner.

\subsection{Purely analytic closure of the exact residual
\texorpdfstring{\(\mathcal D_{20}\)}{D20}}

The final residual is closed without numerical sign tests.  We use a
closed compact superset below \(K=200\) and an analytic curvature
propagation at and above that frequency.  The detailed retained-remainder,
angular-block, scalar-positivity, and base-frequency calculations are
printed in Parts~VII--XIII of the accompanying \emph{Purely Analytic
Supplement}; the logical assembly is included here.

\begin{lemma}[Analytic compact theorem]
\label{cm:lem-compact-certificate}
For
\[
 \rho_c\leq\rho\leq\frac{39}{50},
 \qquad
 k_{11}(\rho)\leq K\leq200,
\]
one has \(N_D(A_{\rho,1},K^2)<W(\rho,K)\).
\end{lemma}

\begin{proof}
Let \(P_{\rm coarse}\) be the strict lattice proxy in
\eqref{cm:eq-coarse}.  Parts~VII--IX of the \emph{Purely Analytic
Supplement} derive an exact retained-remainder identity and a smooth
quantity \(\mathcal H\) such that
\begin{equation}
 \mathcal E_{\rm ang}<\mathcal H
 \quad\Longrightarrow\quad
 P_{\rm coarse}<W.
 \label{cm:eq-analytic-compact-implication}
\end{equation}
Here \(\mathcal E_{\rm ang}\) is the exact error made when the inverse
radial action is sampled on the angular half-integer lattice; the strict
counting convention is built into its definition.

On the displayed closed domain,
\[
 \frac{(1-\rho)k_{11}(\rho)}{\pi}
 =\sqrt{1+\frac{132(1-\rho)^2}{\pi^2}}>1,
\]
so at least one complete radial block is present.  Applying the one-layer
left-block bound across the disjoint complete cells gives the strict global
estimate below; its strict rounding convention includes the common
radial--angular jump faces.
\begin{equation}
 \mathcal E_{\rm ang}
 <\frac{(1-\rho^2)K^2}{6}-\frac12.
 \label{cm:eq-analytic-angular-bound}
\end{equation}
The remaining scalar theorem, proved by monotonicity in \(K\), rational
secant majorants, and positive Bernstein coefficients on the base curve,
gives on the same closed domain
\begin{equation}
 \mathcal H-\frac{(1-\rho^2)K^2}{6}+\frac12>0.
 \label{cm:eq-analytic-scalar-bound}
\end{equation}
Thus \(\mathcal E_{\rm ang}<\mathcal H\), and
\eqref{cm:eq-analytic-compact-implication} yields
\(P_{\rm coarse}<W\).  The phase reduction
\eqref{cm:eq-coarse} proves the claim.  Every inequality used in
\eqref{cm:eq-analytic-angular-bound}--\eqref{cm:eq-analytic-scalar-bound}
is derived in Parts~VII--IX by integration, exact rational comparison,
or positive-side squaring.
\end{proof}

We next state the aggregate argument in the notation needed for its
analytic propagation.  Put
\[
 d_\rho=\frac{2\arcsin\rho}{\pi},\qquad
 \mu=\rho K,\qquad
 b=\frac{2\pi\rho}{1-\rho},\qquad
 \overline R=\mu+\frac12,\qquad
 S=\sqrt{\overline R^2+bK},
\]
and
\[
 \overline{\mathcal I}
 =\frac12\left[
 \overline RS+bK\operatorname{arsinh}
 \frac{\overline R}{\sqrt{bK}}-\overline R^2\right].
\]
For \(K\eta_\rho>1\), define
\begin{equation}
\begin{aligned}
 \mathcal B(\rho,K)={}&
 (\mu-\tfrac12)(K\eta_\rho-1)
 +\frac{d_\rho}{2}(\mu-\tfrac32)(\mu-\tfrac12)\\
 &-(1+d_\rho)\overline{\mathcal I}
 -\frac{8(\mu+\tfrac12)}{15\sqrt\mu}.
\end{aligned}
\label{cm:eq-B-def}
\end{equation}

For completeness, we now prove the discrete bridge behind this reserve.
When \(\mu>1/2\), write uniquely
\[
 \mu-\frac12=J+\tau,\qquad
 J\in\mathbb N_0,\quad0\leq\tau<1,
 \qquad
 R=\begin{cases}J,&\tau=0,\\J+1,&\tau>0,\end{cases}
\]
and put \(x_r=r+1/2\) for \(0\leq r<R\).  Then
\begin{equation}
 n_r:=\lfloor\mu-x_r\rfloor=J-r,\qquad
 q_r:=x_r+n_r=J+\frac12,\qquad
 \sum_{r=0}^{R-1}n_r=\frac{J(J+1)}2,\qquad
 \mu-q_r=\tau.
 \label{cm:eq-aggregate-arithmetic}
\end{equation}
For \(C>0\), set
\begin{equation}
 \mathcal I(C,R)=\frac12\left[
 R\sqrt{R^2+C}+C\operatorname{arsinh}\frac R{\sqrt C}-R^2
 \right].
 \label{cm:eq-I-def}
\end{equation}
With \(C=bK\), define the exact discrete reserve
\begin{equation}
 \mathcal Q=
 R\lfloor K\eta_\rho\rfloor
 +d_\rho\frac{J(J+1)}2
 -(1+d_\rho)\mathcal I(C,R)
 -\frac{8R\tau^{5/2}}{15\sqrt\mu}.
 \label{cm:eq-Q-def}
\end{equation}

\begin{lemma}[Discrete aggregate implication]
\label{cm:lem-Q-implication}
If \(\mathcal Q(\rho,K)\geq0\), then
\begin{equation}
 N_D(A_{\rho,1},K^2)
 \leq P_{\rm coarse}(\rho,K)<W(\rho,K).
 \label{cm:eq-Q-implication}
\end{equation}
\end{lemma}

\begin{proof}
For the accepted first-shelf length \(p_r\) in the low tail beginning at
\(x_r\), equal floors on that shelf give
\[
 1>A_{\rho,K}(x_r)-A_{\rho,K}(x_r+p_r)
   =\int_{x_r}^{x_r+p_r}\vartheta(z)\,dz,
 \qquad \vartheta:=-A_{\rho,K}'.
\]
The fixed-ratio curvature calculation gives
\(\vartheta(0)=0\) and
\(\vartheta'(z)\geq(1-\rho)/(\pi\rho K)=2/C\).  Consequently
\[
 1>\frac{(x_r+p_r)^2-x_r^2}{C},
\]
and hence
\[
 p_r<\sqrt{x_r^2+C}-x_r=:f_C(x_r).
\]
The function \(f_C\) is positive and strictly convex.  The midpoint
inequality therefore yields
\begin{equation}
 \sum_{r=0}^{R-1}p_r
 <\sum_{r=0}^{R-1}f_C(r+1/2)
 \leq\int_0^Rf_C(x)\,dx=\mathcal I(C,R).
 \label{cm:eq-shelf-integral}
\end{equation}
The interface remainder common to all low tails is zero at \(\tau=0\) and
otherwise satisfies
\[
 0\leq\delta_\tau\leq
 \frac{2\tau^{5/2}}{15\sqrt\mu}.
\]
The exact concave-head/convex-tail estimate from the fixed-ratio argument is
\[
 \frac{\mathcal T_r}{2}
 \leq\int_{x_r}^KA_{\rho,K}(z)\,dz+\delta_\tau
 -\frac14\left(
 \lfloor K\eta_\rho\rfloor+d_\rho n_r-(1+d_\rho)p_r
 \right).
\]
This formula also holds when \(n_r=0\), with absent head and \(p_r=0\).
Summing, and then using \eqref{cm:eq-aggregate-arithmetic} and
\eqref{cm:eq-shelf-integral}, gives
\begin{equation}
 \sum_{r<R}\frac{\mathcal T_r}{2}
 -\sum_{r<R}\int_{x_r}^KA_{\rho,K}(z)\,dz
 <-\frac14\mathcal Q(\rho,K).
 \label{cm:eq-low-tail-Q}
\end{equation}

For \(r\geq R\), one has \(x_r\geq\mu\), including equality at a
half-integer interface.  Zero extension and the convex-tail inequality give
\[
 \mathcal T_r\leq2\int_{x_r}^KA_{\rho,K}(z)\,dz;
\]
tails beginning at or above \(K\) vanish.  If
\[
 h_\ell=
 \left[A_{\rho,K}(\ell+\tfrac12)+\tfrac14\right]_<,
\]
the exact multiplicity identity is
\[
 \sum_{\ell\geq0}(2\ell+1)h_\ell
 =\sum_{r\geq0}\left(h_r+2\sum_{\ell>r}h_\ell\right).
\]
Each parenthesis is bounded by \(\mathcal T_r\); at an integer proxy wall
the strict count is one smaller than the ordinary floor.  Hence
\(\mathcal Q\geq0\) and \eqref{cm:eq-low-tail-Q} imply
\[
 P_{\rm coarse}
 <2\sum_{r\geq0}\int_{x_r}^KA_{\rho,K}(z)\,dz
 =2\int_0^Kf_3(z)A_{\rho,K}(z)\,dz.
\]
The weighted scaffold satisfies
\(\int_0^zf_3(t)\,dt\leq z^2/2\).  Since \(A_{\rho,K}\) decreases and
vanishes at \(K\), integration by parts gives
\[
 2\int_0^Kf_3A_{\rho,K}
 \leq\int_0^K2zA_{\rho,K}(z)\,dz=W(\rho,K).
\]
The first inequality is strict, including \(\tau=0\), every half-integer
\(\mu\), and every integer \(K\eta_\rho\).  Finally
\eqref{cm:eq-coarse} gives \(N_D\leq P_{\rm coarse}\).
\end{proof}

If \(K\eta_\rho>1\) and \(\mu>3/2\), then
\[
 R\lfloor K\eta_\rho\rfloor
 >R(K\eta_\rho-1)
 \geq(\mu-\tfrac12)(K\eta_\rho-1),
\]
\[
 \frac{d_\rho J(J+1)}2
 >\frac{d_\rho}{2}(\mu-\tfrac32)(\mu-\tfrac12),\qquad
 \mathcal I(C,R)<\overline{\mathcal I},\qquad
 R\tau^{5/2}<\overline R.
\]
Term-by-term substitution in \eqref{cm:eq-Q-def} therefore proves
\begin{equation}
 \mathcal Q(\rho,K)>\mathcal B(\rho,K).
 \label{cm:eq-B-Q-chain}
\end{equation}
Part~XIII independently develops the continuous curvature propagation
from this analytic discrete bridge.

\begin{lemma}[Purely analytic aggregate tail]
\label{cm:lem-aggregate-tail}
For
\[
 \rho_c\leq\rho\leq\frac{39}{50},
 \qquad K\geq200,
\]
one has \(N_D(A_{\rho,1},K^2)<W(\rho,K)\).
\end{lemma}

\begin{proof}
It is enough to work on the rational superset
\(7/51\leq\rho\leq39/50\), because \(\pi<22/7\) implies
\(7/51<\rho_c\).  On this superset and for \(K\geq200\),
\(\mu>3/2\) and \(K\eta_\rho>1\); the latter follows from
\[
 G_1(s)=\frac1\pi\int_s^1\arccos t\,dt
 \geq\frac{2\sqrt2}{3\pi}(1-s)^{3/2},
 \qquad s\leq\frac{39}{50}.
\]

Differentiating \(\overline{\mathcal I}\) twice gives the exact formulas
recorded in Part~XIII and, since
\(S>\overline R>\rho K\), the curvature estimate
\[
 \mathcal B_{KK}(\rho,K)>F(\rho),
 \qquad
 F(\rho):=
 2\rho\eta_\rho+d_\rho\rho^2
 -\frac{3(1+d_\rho)}{800}\frac{2\pi\rho}{1-\rho}.
\]
Parts~X--XII prove analytically, throughout the rational superset,
\begin{equation}
 F(\rho)>\frac{1549}{84150},\qquad
 \mathcal B_K(\rho,200)>0,\qquad
 \mathcal B(\rho,200)>60.
 \label{cm:eq-analytic-base-signs}
\end{equation}
Those proofs use elementary integral bounds, exact rational
cross-multiplication, positive-side squaring, and explicitly displayed
Bernstein coefficients; no interval enclosure is invoked.

Twice integrating the strict curvature bound from \(200\) to \(K\) yields
\begin{align}
 \mathcal B(\rho,K)
 &=
 \mathcal B(\rho,200)
 +(K-200)\mathcal B_K(\rho,200)
 +\int_{200}^K(K-s)\mathcal B_{KK}(\rho,s)\,ds\notag\\
 &>60+\frac{1549}{168300}(K-200)^2>0.
 \label{cm:eq-analytic-tail-reserve}
\end{align}
This includes \(K=200\) by the strict base value.  Hence
\(\mathcal Q>\mathcal B>0\) by \eqref{cm:eq-B-Q-chain};
\eqref{cm:eq-Q-implication} and \eqref{cm:eq-coarse} finish the proof.
\end{proof}

\begin{proposition}[Exact analytic closure of \(\mathcal D_{20}\)]
\label{cm:prop-exact-closure}
The compact-middle residual is empty after the preceding two lemmas.
\end{proposition}

\begin{proof}
On \(\rho_c\leq\rho<1\),
\[
 z_\rho\geq z_{\rho_c}=\pi+\frac12>\frac72,
 \qquad
 k_{11}(\rho)^2=z_\rho^2+132>\frac{577}{4}>144,
\]
so \(k_{11}(\rho)>12\).  Define the unique owners
\[
 \mathcal C_{21}:=\mathcal D_{20}\cap\{K\leq200\},
 \qquad
 \mathcal T_{21}:=\mathcal D_{20}\cap\{K>200\}.
\]
They form a disjoint partition
\(\mathcal D_{20}=\mathcal C_{21}\mathbin{\dot\cup}\mathcal T_{21}\).
Every point of \(\mathcal C_{21}\) is in the closed domain of
Lemma~\ref{cm:lem-compact-certificate}; every point of
\(\mathcal T_{21}\) is in the domain of
Lemma~\ref{cm:lem-aggregate-tail}.  Thus
\begin{equation}
 \boxed{\mathcal D_{21}:=\mathcal D_{20}\setminus
 (\mathcal C_{21}\cup\mathcal T_{21})=\varnothing.}
 \label{cm:eq-D21-empty}
\end{equation}
The face \(K=200\) belongs only to \(\mathcal C_{21}\).
The inherited faces \(\rho=39/50\), \(K=k_{11}(\rho)\), and
\(K=U(\rho)\) remain outside \(\mathcal D_{20}\), while
\(\rho=\rho_c\) is included subject to its two strict frequency
inequalities.  The exact owner subtraction is proved row by row in
Parts~I--VI of the \emph{Purely Analytic Supplement}.
\end{proof}

\begin{proposition}[Compact-middle P\'olya theorem]
\label{prop:compact-middle}
For every
\[
 \rho_*<\rho<\frac78,\qquad K\geq0,
\]
the Dirichlet counting function of the concentric spherical shell satisfies
\[
 N_D(A_{\rho,1},K^2)\leq W(\rho,K).
\]
\end{proposition}

\begin{proof}
Proposition~\ref{cm:prop-primitive-owners} proves the inequality outside
the exact residual \(\mathcal D_{16}\) in \eqref{cm:eq-D16}.  Within that
residual, Lemmas~\ref{cm:lem-first-bands}, \ref{cm:lem-k5},
\ref{cm:lem-k6}, \ref{cm:lem-lower-d}, \ref{cm:lem-lower-closure},
\ref{cm:lem-k11}, and \ref{cm:lem-optical} remove precisely the regions
subtracted in the definition of \(\mathcal D_{20}\).  Hence any point not
yet owned belongs to
\(\mathcal D_{20}\) in \eqref{cm:eq-D20}.  The analytic compact theorem proves
the inequality on \(\mathcal C_{21}\), and the analytic aggregate-tail lemma
proves it on \(\mathcal T_{21}\).  These two sets are disjoint and exhaust
\(\mathcal D_{20}\), as the exact identity \eqref{cm:eq-D21-empty} shows.
Thus no residual point remains.  The cases on every threshold face have
the explicit owners specified above; in particular, equality at a
frequency wall is not lost in taking complements.  Finally, \(K=0\) is
immediate from \(N_D(A_{\rho,1},0)=W(\rho,0)=0\).  This proves the assertion for
the whole stated parameter range.
\end{proof}

\subsection{Analytic dependency, coverage, and optional cross-checks}
\label{cm:sec-certificate-protocol}

No floating-point computation, interval enclosure, or executable truth value
is a premise of the proof.  The finite portion consists of exact identities,
rational cross-multiplications, positive-side squarings, and explicit sign
decompositions.  Parts~I--VI of the \emph{Purely Analytic Supplement} expose
the former \(611\) ledger predicates in the following normalized disjoint
allocation, after each repeated identity is assigned to one canonical owner:
\[
\begin{array}{c|r}
\text{analytic packet}&\text{uniquely owned rows}\\ \hline
\text{remaining threshold, zero, lower, and cap--74 rows}&103\\
\text{universal arithmetic and checkpoint inventories}&278\\
\text{connected high-event registry}&83\\
\text{seam and two-sided payments}&57\\
\text{upper owner and optical screens}&22\\
\text{exact owner subtraction}&68\\ \hline
\text{total}&611.
\end{array}
\]
The normalized packet census checks that all pairwise intersections are empty and that
the union is the full ledger.  Thus a repeated identity is assigned exactly
once, not counted once per place where it is used.

Parts~VII--IX prove the compact theorem by the exact retained-remainder
identity, the paired angular-block inequality, and scalar positivity on the
closed compact domain.  Parts~X--XIII prove the three base-frequency signs
and propagate them analytically to every \(K\geq200\).  Part~XIV audits the
final disjoint owner split at \(K=200\).  Together these arguments replace,
in the logical dependency graph, the former \(10{,}580\)-box compact
certificate, \(1{,}286\)-box aggregate certificate, and executable
state-classifier.

The earlier exact-arithmetic and interval programs are retained separately
as optional regression tools.  Agreement with them is useful evidence
against transcription errors, but their outputs, precision settings, and
hashes are not mathematical assumptions and are not used in any implication
above.

\section{Proof of the spectral theorem}

We now assemble the regional results without importing any new estimate.
First, elementary bounds give
\[
  3\sqrt3>5>\pi,
\]
and hence \(\omega_0>0\).  Also \(\sqrt3<2\) and \(\pi>3\) give
\(\omega_0<1/6\).  Since \(C_*>0\),
\begin{equation}
  0<\rho_*<\omega_0<\frac16<\frac78.
  \label{eq:ratio-ordering}
\end{equation}
Therefore
\begin{equation}
  (0,1)
  =(0,\rho_*]\mathbin{\dot\cup}
  (\rho_*,7/8)\mathbin{\dot\cup}
  [7/8,1).
  \label{eq:ratio-partition}
\end{equation}

Fix \(0<\rho<1\).  At \(K=0\), positivity in
Proposition~\ref{prop:sc-separated-spectrum} gives
\begin{equation}
  N_D(A_{\rho,1},0)=0=W(\rho,0).
  \label{eq:k-zero}
\end{equation}
Let \(K>0\).  If \(0<\rho\le\rho_*\), apply
Theorem~\ref{thm:sc-small-hole-endpoint}.  If \(\rho_*<\rho<7/8\), apply
Proposition~\ref{prop:compact-middle}.  If \(7/8\le\rho<1\), apply
Lemma~\ref{lem:thin-shell}.  The three cases in
\eqref{eq:ratio-partition} are disjoint and exhaustive, so
\begin{equation}
  N_D(A_{\rho,1},K^2)
  \le\frac{2}{9\pi}(1-\rho^3)K^3
  \qquad(K\ge0).
  \label{eq:unit-shell-theorem}
\end{equation}

Now
\[
  \omega_3=\frac{4\pi}{3},
  \qquad
  L_3=\frac{\omega_3}{(2\pi)^3}=\frac1{6\pi^2},
  \qquad
  |A_{\rho,1}|=\frac{4\pi}{3}(1-\rho^3),
\]
so
\begin{equation}
  L_3|A_{\rho,1}|=\frac{2}{9\pi}(1-\rho^3).
  \label{eq:weyl-normalization}
\end{equation}
Taking \(K=\sqrt\Lambda\) in \eqref{eq:unit-shell-theorem} proves the
unit-shell form of Theorem~\ref{thm:spectral} for every \(\Lambda\ge0\).

For arbitrary \(0<r<R\), put \(\rho=r/R\).  The unitary dilation
\[
  (D_Ru)(x)=R^{-3/2}u(x/R)
\]
maps \(L^2(A_{\rho,1})\) onto \(L^2(A_{r,R})\), and its Dirichlet form
scales by \(R^{-2}\).  Hence
\[
  \lambda_j^D(A_{r,R})=R^{-2}\lambda_j^D(A_{\rho,1}),
\]
and the strict threshold transforms as
\begin{equation}
  N_D(A_{r,R},\Lambda)
  =N_D(A_{\rho,1},R^2\Lambda).
  \label{eq:dilation-count}
\end{equation}
Applying the unit-shell result at \(R^2\Lambda\),
\[
\begin{aligned}
  N_D(A_{r,R},\Lambda)
  &\le\frac{2}{9\pi}(1-\rho^3)(R^2\Lambda)^{3/2}\\
  &=\frac{2}{9\pi}(R^3-r^3)\Lambda^{3/2}\\
  &=L_3|A_{r,R}|\Lambda^{3/2}.
\end{aligned}
\]
This proves Theorem~\ref{thm:spectral} for every \(R>0\), including
\(R<1\).

\section{Proof that spherical shells do not tile}

We prove Theorem~\ref{thm:nontiling} independently of the spectral argument.
Set \(\alpha=r/R\in(0,1)\) and scale by \(1/R\).  It suffices to consider
\[
  T=A_{\alpha,1}=\{x:\alpha<|x|<1\}.
\]
Every rigid motion has the form \(x\mapsto Qx+a\), with \(Q\) orthogonal.
Because \(T\) is radial, \(Q(T)=T\), so every congruent copy is a translate
\[
  T_a=a+T=\{x:\alpha<|x-a|<1\}.
\]

Assume for contradiction that \(\{T_a:a\in\mathfrak C\}\) has
pairwise-disjoint interiors and covers almost everywhere.  Duplicate centers
are impossible because they give the same nonempty open tile.

Fix a unit vector \(e\) and put
\[
  c=\frac{1+\alpha}{2}e,
  \qquad
  \delta=\frac{1-\alpha}{2}.
\]
Write \(B(x,s)=\{y\in\R^3:|y-x|<s\}\) for the open Euclidean ball.
Then \(B(c,\delta)\subset T\).  Therefore \(T_a\) contains
\(B(a+c,\delta)\), and disjointness implies
\begin{equation}
  |a-b|\ge1-\alpha
  \qquad(a\ne b).
  \label{eq:center-separation}
\end{equation}
Thus the centers are locally finite: a bounded set cannot contain infinitely
many points with fixed positive separation.  They are also countable, since
\[
  \mathfrak C
  =\bigcup_{m\ge1}(\mathfrak C\cap\overline B(0,m))
\]
is a countable union of finite sets.

Each tile boundary is a union of two spheres and is three-dimensional null.
The countable union
\[
  Z=\bigcup_{a\in\mathfrak C}\partial T_a
\]
is therefore null.  Exact open-shell coverage is already a special case of
almost-everywhere coverage.  If closed shells cover exactly or almost
everywhere, removing \(Z\) shows that their open interiors still cover almost
everywhere.  Hence it is enough to assume
\begin{equation}
  \left|\R^3\setminus\bigcup_{a\in\mathfrak C}T_a\right|=0.
  \label{eq:ae-coverage}
\end{equation}

Choose one tile \(T_{a_0}\) and its outer sphere
\[
  \Sigma=\{x:|x-a_0|=1\}.
\]
Fix \(p\in\Sigma\) and put \(u=p-a_0\).  For \(n\ge1\), let
\[
  U_n=B\left(p+\frac2n u,\frac1n\right).
\]
Every \(U_n\) lies outside \(T_{a_0}\), is open, and has positive volume.
By \eqref{eq:ae-coverage}, choose \(x_n\in U_n\cap T_{a_n}\) with
\(a_n\ne a_0\).  Since \(x_n\to p\) and \(|x_n-a_n|<1\), the centers
\(a_n\) eventually lie in a fixed bounded set.  Local finiteness lets us
pass to a subsequence with one fixed center \(a\ne a_0\).  Thus
\(p\in\overline{T_a}\).

The point \(p\) is not in \(T_a\).  Otherwise some
\(B(p,\varepsilon)\) lies in \(T_a\).  Choose
\[
  0<t<\min\{\varepsilon,1-\alpha\}.
\]
Then \(p-tu\in T_{a_0}\cap T_a\), contradicting disjoint interiors.
Therefore every \(p\in\Sigma\) belongs to the boundary of a neighboring
tile.

If \(p\in\Sigma\cap\partial T_a\), then
\(|p-a|\in\{\alpha,1\}\), hence \(|a-a_0|\le2\).  Only finitely many
centers satisfy this bound.  Consequently \(\Sigma\) is covered by finitely
many neighboring inner or outer boundary spheres.

Translate so that \(a_0=0\).  The intersection of \(\Sigma\) with a
neighboring sphere \(\{x:|x-a|=s\}\), \(s\in\{\alpha,1\}\), satisfies
\[
  |x|^2=1,
  \qquad |x-a|^2=s^2,
\]
and hence
\begin{equation}
  2a\cdot x=|a|^2+1-s^2.
  \label{eq:sphere-plane}
\end{equation}
Because \(a\ne0\), \eqref{eq:sphere-plane} is a genuine plane.  Its
intersection with \(\Sigma\) is empty, a tangent point, or a circle, and
therefore has zero surface measure on \(\Sigma\).  A finite union of such
sets cannot cover \(\Sigma\), whose area is \(4\pi\).

There is no coincident-sphere exception.  A neighboring outer sphere equals
\(\Sigma\) only when it has the same center, which is the original tile or a
forbidden duplicate.  A neighboring inner sphere cannot equal \(\Sigma\)
because \(\alpha<1\).  We have reached a contradiction.  Scaling back proves
Theorem~\ref{thm:nontiling}.

\appendix

\section{External inputs and verification boundary}
\label{app:external-boundary}

The proof above is self-contained with respect to every shell-specific
spectral reduction, weighted count, endpoint estimate, regional subtraction,
fixed-channel variational comparison, seam assignment, and finite
exact-rational predicate.  It uses the following conventional external results, each only
in the scope stated where it enters the proof:
\begin{enumerate}
  \item standard spherical-harmonic decomposition, regular
        Sturm--Liouville theory, and the classical Bessel identities collected
        in \cite{DLMF};
  \item the uniform one-variable Bessel-phase enclosure of
        \cite{FLPSphase}, together with the explicitly quoted floor-sum and
        turning-point inequalities of \cite{FLPSballs,FLPSannuli};
  \item Lorch's strict lower bound for the first positive zero of
        \(J_\nu\), in precisely the half-integer specializations derived in
        the text \cite{Lorch1993}.
\end{enumerate}

All shell-specific finite decisions are exposed in the accompanying
\emph{Purely Analytic Supplement}.  They require only exact integer and
rational arithmetic, elementary integral inequalities, and the classical
real functions already present in the analytic argument.  The archived
interval programs have no place in this verification boundary; they are
optional regression checks only.

\begin{thebibliography}{5}

\bibitem{FLPSballs}
N.~Filonov, M.~Levitin, I.~Polterovich, and D.~A.~Sher,
\newblock P\'olya's conjecture for Euclidean balls,
\newblock \emph{Invent. Math.} \textbf{234} (2023), 129--169.
\newblock \href{https://doi.org/10.1007/s00222-023-01198-1}{doi:10.1007/s00222-023-01198-1}.

\bibitem{FLPSphase}
N.~Filonov, M.~Levitin, I.~Polterovich, and D.~A.~Sher,
\newblock Uniform enclosures for the phase and zeros of Bessel functions and
their derivatives,
\newblock \emph{SIAM J. Math. Anal.} \textbf{56} (2024), no.~6, 7644--7682.
\newblock \href{https://doi.org/10.1137/24M1642032}{doi:10.1137/24M1642032}.

\bibitem{FLPSannuli}
N.~Filonov, M.~Levitin, I.~Polterovich, and D.~A.~Sher,
\newblock P\'olya's conjecture for Dirichlet eigenvalues of annuli,
\newblock \emph{J. Lond. Math. Soc.} \textbf{113} (2026), no.~2, e70425.
\newblock \href{https://doi.org/10.1112/jlms.70425}{doi:10.1112/jlms.70425}.

\bibitem{Lorch1993}
L.~Lorch,
\newblock Some inequalities for the first positive zeros of Bessel
functions,
\newblock \emph{SIAM J. Math. Anal.} \textbf{24} (1993), no.~3, 814--823.
\newblock \href{https://doi.org/10.1137/0524050}{doi:10.1137/0524050}.

\bibitem{DLMF}
NIST Digital Library of Mathematical Functions,
\newblock Chapter 10: Bessel functions,
\newblock \url{https://dlmf.nist.gov/}, accessed 15 July 2026.

\end{thebibliography}

\end{document}
